% 来源: 19c9fbc3-e932-88ee-8000-0000bb606426_3.2.2_前缀匹配查找.pdf
% 提取时间: 2026-02-28T22:51:49+08:00

3.2.2 前缀匹配查找
考虑伦敦的一个航班数据库，其中列出了飞往 1000个美国城市的航班。一种替代方案是保存一
条记录，指定前往这 1000个城市中每一个城市的路径。假设除了飞往加利福尼亚的航班通过洛杉
矶中转，大多数飞往美国的航班都通过波士顿中转。这一观察可以被利用来将航班数据库从一千条
记录减少到两条前缀记录（USA* → Boston；USA.CA.* → LA）。
这个简化方法的问题在于，像 USA.CA.Fresno这样的目的城市现在会匹配到USA*和USA.CA.*这
两个前缀,此时数据库必须返回最长匹配（USA.CA.*）。因此，前缀可以被用来压缩大型数据库，
但代价是更复杂的最长匹配前缀查找。
正如第 2章中所描述的，互联网使用相同的理念。2022年核心路由器仅存储大约 900,000个前
缀，而不是为每个可能的网络地址存储下数十亿个条目。例如，对于核心路由器来说，加利福尼亚
大学洛杉矶分校（UCLA）这样的大学内的所有计算机，可能都可以通过同一个下一跳来访问。如
果 UCLA 内的所有计算机都被分配了相同的初始位集合（网络号或前缀），那么路由器可以为
UCLA存储一条条目，而不是为 UCLA 的每台计算机存储数千条条目。
自第一版以来，世界发生了显著变化，具体如下：
� 前缀增长 : 自第一版以来，核心路由表的前缀数量从 2004年的大约 150,000个增长到了
2022
年的大约 900,000个。
� IPv6: 在第一版写作时，128位的 IPv6虽然已经被讨论，但由于网络地址转换（NAT）设备的
普及，尚未广泛使用。然而，随着移动设备和物联网（IoT）设备的普及，IPv6的部署速度迅速加
快，到 2022年，Google用户的 IPv6可用率已超过 30%（Google,2022）。
� DRAM与 SRAM:许多路由器硬件查找发现，使用少量片上内存和大容量低延迟的外部DRAM
（动态随机存取存储器）相比SRAM（静态随机存取存储器）更为经济。
� 可编程芯片与 TCAM:例如，英特尔的 Tofino-3（英特尔公司，2022）等芯片已经问世，这
些芯片具有相当大的 TCAM（内容可寻址存储器）并且是可编程的。虽然可以使用 TCAM进行 IP
查找，但新的数据结构可以利用 TCAM扩展到更大的数据库。此外，它们可以使用一种新的用于
路由器编程的高级语言 P4（P4开源编程语言，2022）来编程，执行不同的 IP查找方案。
� 软件定义网络（SDN）:SDN运动允许通过集中控制器改变控制平面（以及通过可编程芯片甚
至数据平面）。每个路由器只需允许一组匹配-动作规则，这些规则可以通过简单地更改匹配的定
义来实现桥接查找、MPLS查找、IP查找或其他形式的查找

                                                     1
� 多核处理器: 处理器变得更快且具有多核。因此，现在在千兆速度下实现 IP查找的软件实现是
可行的。
� 网络功能虚拟化（NFV）: 移动运营商中有一种日益增长的趋势，即用灵活的软件实现来替代
执行网络功能（如最长匹配前缀、防火墙、家长控制等）的复杂中间盒，这种趋势被称为网络功能
虚拟化（NFV，2022）。



尽管发生了这些变化，但基本的算法思想仍然保持不变，只是存在一些我们将指出的小变动，包括
快速的软件实现（如 DXR，2022）和针对大量 DRAM和少量片上 SRAM优化的算法（如 SAIL，
Yang等，2014）。
整个章节的组织如下：第 11.1 节介绍了前缀查找。第 11.2 节描述了减少 IP查找需求的尝试。第
11.3 节展示了基于缓存和并行硬件的非算法技术查找。第 11.4 节描述了基于单比特 Trie 的最简单
技术。
   随后，本章节过渡到描述七种更复杂的方案：多比特 Trie（第 11.5节）、层次压缩Trie（第 11.6
 节）、Lulea压缩 Trie（第 11.7节）、树位图（第 11.8节）、前缀范围上的二分查找（第 11.9节）
                                                      （其现
代表现形式称为 DXR）、前缀长度上的二分查找（第 11.11节）和前缀长度上的线性查找（第 11.12
节）。
本章最后是第 11.13节关于内存分配问题，第 11.14节关于固定功能查找芯片，以及第 11.15节关于
可编程芯片和用于编程的 P4 语言。本章描述的技术（及其相应的原理）在表 11.1 中进行了总
结。




                                                         2
11.1 前缀查找介绍

本节介绍前缀表示法，解释为什么使用前缀查找，并描述评估前缀查找方案的主要指标。



11.1.1 前缀表示法

互联网前缀是用位来定义的，而不是字母数字字符，长度最多为 32 位。然而，令人困惑的是，
IP 前缀通常用点分十进制表示法来书写。例如，在写作时，UCSD（加州大学圣地亚哥分校）有一
个16 位的前缀 132.239。每个点之间的十进制数字代表一个字节。由于二进制中的 132 是
10000100，
239是 11101111，所以 UCSD的前缀在二进制中也可以写作 1000010011101111*，其中星号 *用于
表示其余的位无关紧要。UCSD的主机有以这 16位开头的 32位 IP地址。



                     表 1 本章描述的各种前缀查找方案所涉及的原理
   Table1Principlesinvolvedinthevariousprefix-lookupschemesdescribedinthischapter.

      序号                    原理                             查找方案
     P2a,P10               预计算索引                           标签交换
     P2a,P10         运行时计算传递索引                              IP交换
       P4a           利用 ATM交换机硬件
       P11               缓存整个 IP地址                         查找缓存

                                                                                     3
       P5         硬件并行查找               CAMs
      P4b        扩展前缀以提高速度           受控扩展
       P13        跨度作为自由度    可变跨度 tire
      P4b         压缩以提高速度    Luleatire
     P12P2a       预计算位数集
     P15P12       使用高校查找         前缀长度的二分查找
      P2a      添加标记状态预计算标记观察
      P2a       预计算范围到前缀匹配          前缀上的二分查找

由于前缀可以是可变长度的，表示前缀的另一种常见方法是使用斜杠表示法，其形式为A/L。在这
种情况下，A表示点分十进制表示法的 32位 IP地址，L表示前缀的长度。因此，UCSD的前缀也可
以表示为 132.239.0.0/16，其中长度 16表示只有前 16位（即 132.239）是相关的。描述前缀的第
三种常见方法是使用掩码代替显式的前缀长度。因此，UCSD的前缀也可以描述为 132.239.0.0，
掩码为 255.255.0.0。由于 255.255.0.0的前 16位是 1，这隐含地表示前缀长度为 16位。
在这三种表示前缀的方法中（二进制后面加通配符、斜杠表示法和掩码表示法），后两种在记录大
型前缀时更为紧凑。然而，出于教学原因，使用小型前缀作为示例并以二进制形式书写要容易得
多。因此，在本章中，我们将使用 01110*来表示匹配所有以 01110开头的 32位 IP地址的前缀。读
者应该能够轻松地将这种表示法转换为供应商使用的斜杠或掩码表示法。此外，请注意，大多数前
缀至少有 8 位长；然而，为了使我们的示例更简单，本章使用较小的前缀。11.1.1 为什么前缀
是变长的？
在考虑如何处理可变长度前缀匹配的复杂性之前，值得理解为什么互联网前缀是可变长度的。给定一个
电话号码，例如 858-549-3816，提取前三位数字（即 858）作为区号是很简单的。如果固定长度的前
缀更容易实现，那么可变长度前缀的优势是什么呢？
这个问题的一般答案是，可变长度前缀可以更有效地利用地址空间。这是因为具有大量终端的区域可以
分配较短的前缀，而只有少量终端的区域可以分配较长的前缀。
具体的答案来自互联网地址的历史。互联网起初采用一个简单的层次结构，其中 32 位地址被分为网络地
址和主机号；路由器只存储网络的条目。为了实现灵活的地址分配，网络地址采用了可变大小：A 类
（8 位）、B 类（16 位）和 C 类（24 位）。为了应对 B 类地址的枯竭，无类别域间路由
（CIDR）方案（Rekhter 和 Li，1996）为新的组织分配了多个连续的 C 类地址，这些地址可以通过一
个公共前缀进行聚合，从而减少了核心路由器表的大小。

                                                         4
如今，地址空间可能枯竭的潜在风险使得互联网注册机构在分配 IP 地址时非常保守。一个小型组
织可能只会分配到一小部分 C 类地址，例如一个/30 前缀，这仅允许组织内有四个 IP 地址。许多
组织通过使用网络地址转换（NAT）等方案，在多台计算机之间共享少量 IP 地址来应对这些稀少
的分配。
因此，CIDR 和 NAT 帮助互联网在有限的 32 位地址空间内应对了指数级的增长。虽然采用具有
128 位地址的 IP（IPv6）的速度正在迅速增加，但在短期内 NAT 的有效性以及推出新协议的复杂
性，使得在撰写本文时，32 位 IP 地址仍然占据主导地位。
结论是，部署 CIDR 的决定帮助挽救了互联网，但它也引入了最长匹配前缀查找的复杂性。
11.1.2 查找模型
回顾第 2 章中的路由器模型。一个数据包到达输入链路。每个数据包携带一个 32 位的互联网
（IP）地址。处理器查询转发表以确定数据包的输出链路。转发表包含一组前缀及其对应的输出链
路。数据包与与数据包中目的地址最长匹配的前缀进行匹配，并将数据包转发到相应的输出链路。
确定输出链路的任务称为“地址查找”，这是本章的主题，本章将综述查找算法，并证明查找可以
在千兆和太比特速度下实现。
在寻找 IP 查找解决方案之前，了解一些关于流量分布、内存趋势和数据库大小的基本观察结果是
很重要的，这些内容在表 11.2 中显示。这些观察结果将促使我们了解查找方案的需求。
表 2 关于查找问题的一些当前的数据及其对应解决方案的提示
Table 2 Some current data about the lookup problem and the corresponding implications
for lookup solutions.
观察结果 对查找解决方案的推论
1. 骨干网中有 25 万个并发流
2. 503. 查找受内存访问控制
4. 前缀长度从 8 到 32
5. 今天有 100 万个前缀和多播及主机路由
6. 不稳定的 BGP，多播
7. 更高的速度需要 SRAM
8. IPv6，多播延迟 缓存在骨干路由器中效果并不好查找速度由内存访问次数衡量
简单方案需要进行 24 次内存访问
随着增长，需要 50 万到 100 万个前缀更新时间在毫杪到秒之间
值得最小化内存
今天 32 位和 128 位查找都至关重要
首先，早在 1997 年进行的一项骨干网流量研究（Thompson 等，1997）显示，大约有 25 万个
短时并发流，使用的是相当保守的流量测量方法。测量数据显示，这一数字正在不断增加，轻松超
过了 100 万个并发 TCP 流。这么多的流量意味着缓存解决方案效果不好。
                                                                                   5
其次，同一研究（Thompson 等，1997）显示，路由器接收到的大约一半的数据包是最小尺寸的
TCP 确认。因此，路由器有可能接收到一串最小尺寸的数据包。因此，能够在转发最小尺寸数据
包的时间内进行前缀查找，可以避免输入链路队列的需求，从而简化系统设计。第二个原因是简单
的市场营销：许多供应商声称支持线速转发，这些声称是可以测试的。假设线速转发，以 10
Gbps
（OC-192 速度）转发一个 40 字节的数据包应不超过 32 纳秒，以 40 Gbps（OC-768）转发则
应不超
过 8 纳秒。
显然，对于查找方案来说，最关键的指标是查找速度。第三个观察指出，由于当今计算的成本

主要由内存访问决定，查找速度的最简单衡量标准是最坏情况下的内存访问次数。第四个观察显
示，骨干网数据库中的前缀长度从 8 到 32 不等，因此，天真的方案在最坏情况下需要进行 24 次
内存访问，以尝试所有可能的前缀长度。
第五个观察指出，虽然当前数据库大约有 900,000 个前缀，但可能使用主机路由（完整的 32 位
地址）和多播路由，这意味着未来的骨干路由器将拥有超过 100 万个前缀的前缀数据库。
第六个观察指的是查找数据结构的更新速度，例如添加或删除前缀。不稳定的路由协议实现可能会
导致对更新速度有毫秒级的要求。请注意，无论是秒级还是毫秒级，这都比查找要求低了几个数量
级，这使得实现可以奢侈地在数据结构中预计算（P2a）信息以加速查找，但代价是较长的更新时
间。
第七个观察来自第 2 章。虽然标准的 DRAM 内存便宜，但 DRAM 的访问时间目前约为 20– 40
纳秒，因此在更高速度下可能需要更高速的内存（例如，片上或片外 SRAM，1–5 纳秒）。尽管
DRAM 内存基本上是无限的，但 SRAM 和片上内存由于成本或不可用性而受到限制。因此，第三
个指标是内存使用，其中内存可以是昂贵的高速内存（软件中的缓存，硬件中的 SRAM），也可
以是便宜的慢速内存（例如，DRAM，SDRAM）。
请注意，不进行增量更新的查找方案将需要查找数据库的两个副本，以便在一个副本上进行搜索的
同时在另一个副本上进行查找。因此，仅仅为了将高速内存减少一半，进行增量更新可能是值得
的！
第八个观察涉及前缀长度。IPv6 需要 128 位前缀。多播查找需要 64 位查找，因为完整的组地址
和源地址可以连接成一个 64 位前缀。然而，IPv6 和多播的全面部署仍在进行中。因此，在撰写
本文时，能够支持 32 位和 128 位 IP 查找是很重要的。
总结起来，有趣的指标按重要性排序为查找速度、内存和更新时间。举一个具体的例子，一个好的
片上设计使用 16 兆位的片上内存，可能支持任意一组 1,000,000 个前缀，能够在 8 纳秒内完成
查找，从而在太比特速度下提供线速转发，并且允许前缀在 1 毫秒内更新。
在报告 IP 查找算法的理论性能时，使用以下符号：N 表示前缀的数量（例如，2022 年大型数据

                                                       6
库中的 900,000 个前缀），W 表示地址的长度（例如，IPv4 为 32）。
最后，还有两个额外的观察可以被利用来优化期望情况。
  1. 几乎所有前缀都是 24 位或更少，其中大多数是 24 位前缀，接下来最大的峰值是 16 位。
    一些
    供应商使用这点来显示只有 24 位前缀的最坏情况查找时间；然而，将来可能会导致具有大量
    主机路由（32 位地址）和集成 ARP 缓存的数据库。
  2. 前缀为其他前缀的前缀（例如 00* 和 0001*）的情况相当罕见。事实上，当前数据库中某
    个前缀的最大前缀数是七。
虽然理想的方案是满足最坏情况下查找时间要求，但也希望能够利用这些观察来提高平均存储性能
的方案。
11.1 优化查找
系统设计人员的第一个本能不是解决复杂问题（如最长匹配前缀），而是消除问题。
注意，在像ATM 这样的虚拟电路网络中，当一个源希望向一个目的地发送数据时，会建立一个类
似于电话呼叫的连接。每个路由器上的呼叫编号虚拟电路索引（VCI）是一个中等大小的整数，查
找起来很容易。然而，这需要付出在发送数据之前进行呼叫建立的往返延迟的代价。
根据我们的原则，ATM网络的前一跳交换机将索引（P10，在协议头中传递提示）传递到下一跳交
换机。该索引在前一跳交换机发送数据之前预先计算（P2a）（P3c，将计算转移到空间中）。同
样的抽象理念可以用于数据报网络（如互联网）来简化前缀查找的需求。我们现在描述这个抽象理
念的两个实例：标签交换（第 11.2.1节）和流交换（第 11.2.2节）。




图 1通过让每个路由器将一个索引传递到下一个路由器的转发表中，来取代数据报路由器中对目的
地查找的需求

                                                  7
Figure 1Replacing the need for a destination lookup in a datagram router by having each
router pass an index into the next router’s forwarding table.

11.1.1 线程索引和标签交换
在线程索引（Chandranmenon和 Varghese，1996）中，每个路由器将一个索引传递到下一个路
由器的转发表，从而避免前缀查找。索引由路由协议在拓扑变化时预先计算。因此，在图 11.1中，
源 S向目的地 D发送数据包到第一个路由器 A时，数据包头还包含了一个指向 A的转发表的索引
i。A的 D条目显示下一跳是路由器 B，B在其转发表中存储的 D条目位于索引 j处。因此，A将数据
包发送到 B，但首先将 j写为数据包索引（图 11.1）。这个过程重复进行，路径上的每个路由器都
使用数据包中的索引查找其转发表。
线程索引和虚拟电路索引（VCIs）之间的两个主要区别如下。首先，线程索引是针对每个目的地
的，而不像ATM等虚拟电路网络中那样针对每个活跃的源-目的地对。其次，更重要的是，线程



索引在拓扑变化时由路由协议预先计算。例如，考虑图 11.2，它展示了一个样本路由器拓扑，其中
路由器运行 Bellman–Ford 协议来找到它们到达目的地的距离。
在 Bellman–Ford 算法中（例如，在域内协议路由信息协议 [RIP] 中使用，Perlman，1992），
路由器 R 通过取经过每个邻居到达 D 的成本的最小值来计算其到 D 的最短路径。通过一个邻居
（如 A）的成本是 A到 D的成本（即 5）加上从 R到 A的成本（即 3）。在图 11.2中，R到 D的最
低成本路径是通过路由器 B，成本为 7。R 能够计算这个，因为 R 的每个邻居（如 A，B）将其当
前到 D 的成本传递给 R，如图所示。为了计算索引，我们修改基本协议，使每个邻居除了报告到
目的地的成本外，还报告其到目的地的索引。因此，在图 11.2中，A传递 i，B传递 j；因此，当 R
选择 B 时，它还在其路由表条目中使用 B 的索引 j 来表示 D。总之，每个路由器使用最小成本邻
居的索引作为每个目的地的线程索引。




                                                                                      8
             图 2 通过修改 Bellman–Ford路由来设置线程索引或标签。
    Figure2SettingupthethreadedindexesortagsbymodifyingBellman–Fordrouting.


思科后来引入了标签交换（Meiners等，2008），其概念与线程索引相似，不同之处在于标签交
换还允许路由器传递多个下游路由器的标签（索引）堆栈。然而，这两种方案都不能很好地处理层
次结构。考虑一个从骨干网到达出口域中第一个路由器的数据包。出口域是数据包穿越的最后一个
自主管理网络——假设是数据包目的地所在的企业网络。
   在出口域中的第一个路由器R处避免查找的唯一方法是让出口域外的某个早期路由器将索引
（针对目标子网）传递给R。但这是不可能的，因为先前的骨干路由器应该只有一个针对整个目标
                                              域



的聚合路由条目，因此只能为该域内的所有子网传递一个索引。唯一的解决方案是为域外的路由器
添加额外条目（不可行）或在域入口点进行普通的 IP查找（被选择的解决方案）。今天，标签交
换在一种更通用的形式中蓬勃发展，称为多协议标签交换（MPLS）（Meiners等，2008）。然
而，无论是标签交换还是 MPLS，都无法完全避免普通 IP 查找的需求。

11.1.2 流交换

                                                                              9
第二个优化查找的提案被称为流交换（Newman 等，1997；Parulkar 等，1995）。流交换也依赖
前一跳路由器将一个索引传递到下一跳路由器。然而，与标签交换不同，这些索引是在数据到达时
按需计算的，然后进行缓存。
流交换从包含内部 ATM交换机和（可能较慢的）能够进行 IP转发和路由的处理器的路由器开始。
图 11.3中展示了两个这样的路由器 R1和 R2。当 R2首次向目的地 D发送 IP数据包并到达 R1 的左
输入端口时，输入端口将数据包发送到其中央处理器。这是慢速路径。处理器进行 IP 查找，并在
内部将数据包切换到输出链路 L。到目前为止，这一切都很正常。




图 3在 IP交换中，如果 R1希望切换发送到 D并且目标为输出链路 L的数据包，R1会选择一个空闲
的虚拟电路 I，将映射 I,L 放置在其输入端口，然后将 I 返回给 R2。如果 R2 现在发送标有 VCI I
的数据包到 D，这些数据包将直接被切换到输出链路，而无需经过处理器
Figure3In IP switching, if R1 wishes to switch packets sent to D that are destined for
output link L, R1 picks an idle virtual circuit I, places the mapping I,L in its input port, and
then sends I back to R2.If R2 now sends packets to D labeled with VCI I,the packet will get
switched directly to the output link without going through the processor.

如果 R1决定“切换”发往 D的数据包，情况会变得更加有趣。例如，如果有大量流量发往 D， R1可
能会决定这样做。在这种情况下，R1首先选择一个空闲的虚拟电路标识符 I，将映射 I → L放置在
其输入端口硬件中，然后将 I发送回 R2。如果 R2现在将标有 VCII的数据包发送到 R1的输入端
口，输入端口将查找从 I 到 L 的映射，并将数据包直接切换到输出链路 L，而无需经过处理



器。

                                                                                             10
当然，R2可以与路径中的前一个路由器重复这一切换过程，依此类推。最终，在一系列流交换路
由器的切换部分，IP 转发可以被完全省略。
尽管流交换的设计很优雅，但在骨干网中可能效果不佳。这是因为骨干网流量的寿命较短且呈现出
较差的局部性。相反的观点在 Molinero-Fernandez 和 McKeown（2002）中提出，作者主张基
于 TCP连接的流交换的复兴。他们声称，目前使用电路交换光开关连接核心路由器、骨干网链路
仅以 10% 的容量运行的低利用率以及不断增加的光带宽，都支持在更高速度下电路交换的简单
性。
IP交换和标签交换都是通过在协议头中传递信息来优化 IP查找需求的技术。与 ATM类似，这两种
方案都依赖于传递索引（P10）。然而，标签交换在较早的时间尺度（拓扑变化时间）上预先计算
索引（P2a），而 ATM是在数据传输之前才进行计算。另一方面，在 IP交换中，索引是在数据开
始流动后按需计算的（P2c，惰性求值）。然而，无论是标签交换还是 IP交换，都不能完全避免
前缀查找，并且每种方案都增加了复杂的协议。现在我们重新审视 IP查找的复杂性问题。

11.1.3 标签交换、流交换和多协议标签交换的状态
虽然标签交换和 IP交换最初是为了加快查找速度而引入的，但 IP交换已经消亡。然而，标签交换
以更通用的多协议标签交换（MPLS）（IETFMPLS章程，1997）的形式重新发明了自己，作为提
供流量区分以提供服务质量的机制。正如 VCI提供了一个简单的标签以快速区分流量一样，标签允
许路由器轻松隔离流量以提供特殊服务。实际上，MPLS使用标签来优化数据包分类的需求（第
12章），这比前缀查找要难得多。因此，尽管仍然需要进行前缀匹配，但 MPLS在今天的核心路
由器中也是必不可少的。
简而言之，所需的 MPLS快速路径转发过程如下：识别具有 MPLS头的数据包，提取 20位标签，
并在映射标签到转发规则的表中查找该标签。转发规则指定下一跳，还指定在 MPLS数据包的当
前标签集合上执行的操作。这些操作可以包括移除标签（“弹出标签栈”）或添加标签（“压入标签
栈”）。
路由器的 MPLS实现必须对这一通用过程施加一些限制，以保证线速转发。可能的限制包括要求
标签空间密集，支持的标签数量少于 220（这允许使用更少的查找内存，同时避免使用哈希表），
以及限制在单个数据包上可以执行的标签堆叠操作的数量。



11.2 前缀匹配的非算法技术

                                                       11
在本节中，我们考虑两种不依赖算法方法的前缀查找系统技术：缓存和三态 CAMs。缓存依赖于
地址引用的局部性，而 CAMs 依赖于硬件并行性。

11.2.1 缓存
可以通过使用将 32 位地址映射到下一跳的缓存（P11a）来加速查找。然而，由于骨干网中流量缺
乏局部性，缓存命中率较低（Newman等，1997）。使用大型缓存仍然需要使用精确匹配算法进



行查找。一些研究者提倡为此目的巧妙地修改 CPU缓存查找算法（Chiueh和 Pradhan，1999）。
总的来说，缓存可以有所帮助，但并不能避免快速前缀查找的需求。

11.2.2 三态内容可寻址存储器（CAMs）
三态内容可寻址存储器（TCAMs）允许使用“无关”位，在一次内存访问中提供并行搜索。如今的
TCAMs可以在一个内存周期内进行搜索和更新（例如，10纳秒）并处理任意组合的 100,000个前
缀。它们甚至可以级联形成更大的数据库。然而，TCAMs存在以下问题。




图 4示例前缀数据库用于本章的其余部分。请注意，为了清晰起见，已省略每个前缀对应的下一
跳。Figure4
Sampleprefixdatabaseusedfortherestofthischapter.Notethatthenexthopscorrespondingtoeac
hprefixhave been omitted for clarity.

� 密度扩展：一个 TCAM位需要 10-12个晶体管，而 SRAM需要 4-6个晶体管。因此，TCAM
的密度比 SRAM低，或者占用更多的面积。对于许多路由器来说，板面积是一个关键问题。
� 功耗扩展：由于并行比较，TCAM的功耗较高。然而，CAM供应商正在通过找到关闭部分
CAM

                                                                                    12
的方法来减少功耗，以解决这一问题。功耗是大型核心路由器中的一个关键问题。
� 匹配仲裁：CAM中的匹配逻辑要求所有匹配规则进行仲裁，以确保最高匹配胜出。老一代的
CAM执行一次操作大约需要10纳秒，但对于现代TCAM来说，这已经不再是一个大问题。
� 额外芯片：鉴于许多路由器，如思科 GSR和瞻博 M160，已经有一个专用的应用程序特定集
成电路（ASIC）或网络处理器进行数据包转发，集成分类算法与查找算法，而不增加 CAM 接口
和 CAM 芯片是很有吸引力的。需要注意的是，CAM 通常除了基本的 CAM 芯片外，还需要一个
桥接 ASIC，有时还需要多个 CAM 芯片。
� 内置TCAM的可编程芯片：与上一点提到的额外芯片问题（由分离的 CAM和数据包转发芯片
引起）相比，近年来出现了能够使用内置 TCAM执行转发的可编程芯片，这是一个重大的变化。
例如，英特尔的 Tofino-3（英特尔公司，2022）包含 384个大小为 44×512的 TCAM块，足以
支持大型企业或数据中心，但没有一些算法技巧的话，无法满足骨干网的需求。
总之，CAM技术正在迅速改进，并且在较小的路由器中逐渐取代算法方法。然而，对于可能希望
在未来拥有百万路由数据库的大型核心路由器来说，最好采用能够与 SRAM 等标准内存技术相适
应的解决方案（如我们在本章中所描述的）。SRAM可能始终比 CAM更便宜、更快和更密集。虽



然现在预测算法方法和 TCAM 方法之间的竞争结果还为时过早，但即使是半导体制造商也在两者
之间押注，提供基于算法和基于 CAM 的解决方案。




                                                     13
11.3 单比特 Trie




对于前缀查找的算法技术（P15）调查，最简单的技术是单比特 Trie。考虑图 11.4中的示例前缀数
据库。这个数据库将用于说明本章中的许多算法解决方案。它包含九个前缀，称为 P1到 P9，图中
显示了位串。



实际上，每个前缀都有一个与之关联的下一跳，但在图中省略了。为了避免混乱，前缀名称用来表
示下一跳。因此，在图中，地址 D以 1开头，后面跟着 31个零，将匹配 P4、P6、P7和 P8。最长
匹配是 P7。



                                                 14
图 11.5展示了图 11.4示例数据库的单比特 Trie。单比特 Trie基于分而治之的简单算法技术
（P15），按目的地地址的位进行划分，从最高有效位开始。单比特 Trie 是一个树，其中每个节
点是一个包含 0指针和 1指针的数组。在根节点，所有以 0开头的前缀存储在由 0指针指向的子
Trie中，所有以 1 开头的前缀存储在由 1 指针指向的子 Trie 中。




图 5 图中的示例数据库的单比特 Trie如下所示。
Figure5 Theone-bittrieforthesampledatabaseofFig


       然后，每个子Trie使用分配给该子Trie的前缀的剩余位以类似的方式递归构建。例如，在图
   11.4中，注意到 P1 = 101存储在一条路径中，该路径通过根节点的 1指针、根节点右子节点的 0
指针和路径中下一个节点的 1指针进行跟踪。
需要注意的还有两个细节。在某些情况下，一个前缀可能是另一个前缀的子串。例如，P4 = 1*是
P2 = 111*的子串。在这种情况下，较短的字符串 P4存储在通往较长字符串路径上的一个 Trie节点
中。例如，P4存储在根节点的右子节点中；注意，到这个右子节点的路径是字符串 1，这与 P4相
同。
最后，对于像 P3 = 11001这样的前缀，在我们跟随前三个位后，可能会天真地期望找到对应最后
两位的节点字符串。然而，由于没有其他前缀与 P3共享超过前三个位，这些节点每个只会包含一
                                                    15
个指针。这样的 Trie 节点字符串，每个只有一个指针，称为单向分支。
为了消除这种明显的浪费（P1），可以使用一种简单的技术来压缩单向分支。
在图 11.5中，通过使用一个文本字符串（即“01”）来表示单向分支中的指针，从而实现压缩。因
此，在图 11.5 中，路径上通往 P3 的两个 Trie 节点（每个节点包含两个指针）被一个 2 位的文



本字符串所替代。显然，通过这种转换没有丢失任何信息。（作为练习，确定 Trie 中是否有其他
路径可以类似地压缩。）
要查找目标地址 D的最长匹配前缀，使用 D的位来跟踪 Trie中的路径。路径从根节点开始，继续
直到搜索失败，搜索失败的情况包括到达一个空指针或文本字符串不完全匹配。在跟踪路径的过程
中，算法记录在路径节点上遇到的最后一个前缀。当搜索失败时，这就是返回的最长匹配前缀。
例如，如果 D 以 1110 开头，算法首先跟随根节点的 1 指针到达包含 P4 的节点。算法记住 P4，
并使用 D 的下一个位（1）跟随 1 指针到下一个节点。在这个节点上，算法跟随 1 指针到达 P2。
当算法到达 P2时，用新找到的前缀（P2）覆盖先前存储的值（P4）。此时搜索终止，因为 P2没
有外部指针。
另一方面，考虑搜索一个目标地址 D，其前 5位是 11000。再次使用第一个 1位到达包含 P4
的节点。记住P4作为最后遇到的前缀，然后跟随 1指针到达高度为 2的最右节点。
算法现在跟随 D的第三位（0）到达包含“01”的文本字符串节点。因此，记住 P9作为最后遇到的
前缀。D的第四位是 0，与“01”的第一位匹配。然而，D的第五位是 0（而不是“01”的第二位中的
1）。因此，搜索终止，P9为最长匹配前缀。
Trie的文献（Knuth,1973）没有使用文本字符串来压缩单向分支，如图 11.5所示。相反，经典方案
称为 Patricia trie，使用跳过计数。该计数记录相应文本字符串中的位数，而不是位本身。例如，
我们示例中的文本字符串节点“01”在 Patriciatrie 中将被跳过计数“2”替代。
只要Patricia trie 用于精确匹配，这是完全可行的，这是它们最初的用途。当搜索到达跳过计数节
点时，它会跳过适当数量的位，并跟随跳过计数节点的指针继续搜索。由于跳过的位没有进行匹配
比较，Patricia 要求在搜索结束时，对搜索到的键和 Patricia 找到的条目进行完整的比较。
不幸的是，这在前缀匹配中效果非常差，而前缀匹配并不是 Patriciatries最初设计要处理的应用。
例如，在图 11.5的 Patricia等价物中搜索前 5位为 11000的 D，搜索会跳过最后两位并到达 P3。
此时比较会发现 P3 不匹配 D。

                                                       16
当这种情况发生时，在 Patricia trie 中搜索必须回溯并返回 Trie 向上搜索可能的较短匹配。在这
个例子中，似乎搜索可以记住 P4。但如果 P4也在包含跳过计数节点的路径上遇到，算法甚至不
能确定 P4。因此，必须回溯以检查 P4 是否正确。
不幸的是，BSD的 IP转发实现（Wright和 Stevens，1995）决定使用 Patriciatries作为最佳匹配
前缀的基础。因此，BSD 实现使用了跳过计数；实现还通过用零填充前缀来存储它们。前缀也存
储在 Trie的叶节点中，而不是节点内，如图 11.5所示。结果是，前缀匹配在最坏情况下可能导致
在 Trie 中回溯，最坏情况下需要 64 次内存访问（32 次向下和 32 次向上）。
鉴于使用文本字符串避免回溯的简单替代方案，使用跳过计数是一个糟糕的主意。本质上，这是因
为跳过计数转换没有保留信息，而文本字符串转换保留了信息。然而，由于 BSD的巨大影响，许
多供应商甚至其他算法（例如，Ref.Nilsson和 Karlsson，1998）在其实现中使用了跳过计数。




11.5 多比特 Tire


大多数大容量存储器使用 DRAM。与寄存器访问时间（2纳秒）相比，DRAM的延迟较大（大约
30 纳秒）。由于单比特 Trie 可能需要进行 32 次访问以匹配 32 位前缀，单比特 Trie 的最坏情
况查找时间至少为 32 * 30 = 0.96 微秒。这显然促使了多比特 Trie 查找的需求。一次可以查找
的位数是算法技术（P13）可以利用的一个自由度。
为了以 4 位为步长查找 Trie，主要问题是处理像 10101*（长度为 5）这样的前缀，其长度不是选
择的步长 4 的倍数。如果我们每次查找 4 位，如何确保不会遗漏像 10101* 这样的前缀？受控前
缀扩展解决了这个问题，通过将现有前缀数据库转换为具有更少前缀长度但可能有更多前缀的新数
据库。通过消除所有不是所选步长倍数的长度，扩展允许更快的多比特 Trie 查找，代价是数据库
大小增加。

例如，移除奇数前缀长度将前缀长度从 32减少到 16，并允许每次查找 2位。要移除长度为 3的前
缀如 101*，注意 101* 表示以 101 开头的地址，这又表示以 1010* 或 1011* 开头的地址。因此，

                                                            17
长度为 3的 101*可以被两个长度为 4的前缀（1010*和 1011*）替代，两者都继承了 101*的下一跳
转发条目。

然而，扩展的前缀可能会与新长度的现有前缀冲突。在这种情况下，扩展前缀将被移除。现有前缀
优先，因为它原本长度更长。

本质上，扩展用内存换时间（P4b）。相同的思想可以用于移除除长度 32之外的任何选定长度。
由于 Trie 查找速度线性依赖于长度数量，扩展减少了查找时间。
考虑图 11.4所示的示例前缀数据库，其中包含九个前缀，P1到 P9。图 11.6左侧重复了相同的数据
库。图 11.6 右侧的数据库是通过扩展原始数据库只包含长度为 3 和 6 的前缀构建的等效数据库。
注意 P6 = 1000*扩展到 6位的四个扩展中，有一个与 P7 = 100000*冲突，因此被移除。




11.5.1 固定步长 Trie


    图 11.7展示了与图 11.6相同数据库的 Trie，使用了步长为 3的固定步长 Trie。因此，每个 Trie
   节点使用 3个位。图 11.7中Trie节点内的重复条目与图 11.6右侧的扩展前缀完全对应。例如，图
11.6中的 P6有三个扩展（100001、100010、100011）。




                                                           18
图 6对左侧显示的原始前缀数据库进行受控扩展（其前缀长度为 1、3、4、5和 6）得到扩展数据
库（其前缀长度仅为 3和 6）。
Figure6Controlled expansion of the original prefix database shown on the left (which has
five prefix lengths, 1, 3, 4, 5, and 6)to an expanded database (which has only 2 prefix
lengths, 3 and 6).




这些三个扩展前缀由图 11.7 根节点中的 100 指针指向（因为这三个扩展前缀都以 100 开头），并
存储在根节点右子节点的 001、010和 011条目中。还请注意，根节点中的 100条目存储了一个前
缀 P8（除了指向 P6扩展的指针外），因为 P8 = 100*本身就是一个扩展前缀。

因此，每个 Trie节点元素是包含两个条目的记录：一个存储的前缀和一个指针。Trie查找每次处
理 3 个位。每次跟随指针时，算法都会记住存储的前缀（如果有的话）。当查找在空指针处终止
时，返回路径中的最后一个存储前缀。


                                                                                       19
例如，如果地址 D 以 1110 开头，查找 D 从根节点的 111 条目开始，该条目没有外部指针，但有
一个存储前缀（P2）。因此，D 的查找以 P2 终止。查找以 100000 开头的地址会跟随根节点中
的 100 指针（并记住 P8）。这会导致到达右下角的节点，其中 000 条目没有外部指针，但有一
个存储前缀（P7）。查找以结果 P7终止。可以使用宽内存一次内存访问同时检索指针和存储前缀
（P5b）。

固定步长 Trie的一个特例由 Gupta等人（1998）描述，使用 24、4和 4的固定步长。作者观察
到，具有超过 224位置的 DRAM正在出现，使得 24位步长也是可行的。



11.5.2 可变步长 Trie




              图 7扩展 Trie（每个步长为 3位）对应于图 11.6中的前缀数据库。
Figure7
Expandedtrie(whichhastwostridesof3bitseach)correspondingtotheprefixdatabaseofFig.11.6.

在图 11.7中，最左边的叶节点需要存储 P3 = 11001*的扩展，而最右边的叶节点需要存储 P6
（1000*）和 P7（100000*）。因此，由于 P7，最右边的叶节点需要检查 3 位。然而，最左边的
叶节点没有理由检查超过 2位，因为 P3只有 5位，且根步长为 3位。这里有一个可以优化的自由
度
（P13）。

                                                                                         20
在可变步长Trie中，每个Trie节点检查的位数可以不同，即使是同一层的节点也是如此。为此，
Trie 节点的步长与指向该节点的指针一起编码。可以通过用一个四元素数组替换最左边的节点，
并将步长 2 与指向最左边节点的指针一起编码，将图 11.7 转换为可变步长 Trie（图 11.8）。步长
编码花费 5 位。然而，图 11.8 的可变步长 Trie 比图 11.7 的 Trie 少了四个数组条目。
我们的示例激发了选择步长以最小化 Trie 内存总量的问题。由于扩展用内存换时间，为什么不通
过优化一个自由度（P13），即每个节点使用的步长，来最小化所需的内存呢？为了选择可变步
长，设计者首先要指定最坏情况下的内存访问次数。例如，对于 1 Gbps 的 40 字节数据包和 80
纳秒的
DRAM，我们有 320 纳秒的时间预算，这只允许四次内存访问。这限制了任何搜索路径中的最大
节点数（在我们的示例中为四个）。
在给定这个固定高度的情况下，可以选择步长以最小化存储。这可以通过动态规划（Srinivasan
和 Varghese，1999）在几秒钟内完成，即使对于包含 150,000个前缀的大型数据库也是如此。通
过优化自由度（步长）来最小化在给定最坏情况下树高度的内存使用量。
一个 Trie被称为高度为 h和数据库 D的最优 Trie，如果它在所有针对数据库 D、高度不超过



h 的可变步长 Trie 中具有最小的存储量。很容易证明（见习题）图 11.8 的 Trie 对于图 11.6 左侧
的数据库和高度 2 是最优的。


选择最优步长 Trie 的通用问题可以递归解决（图 11.9）。假设树的高度必须为 h。算法首先选择
步长为 s的根节点。根节点中的 y = 2s个可能指针可以指向 y个非空子 TrieT1,...Ty。如果 s位指
针 pi指向子 TrieTi，那么所有以 pi开头的原始数据库 D中的前缀都必须存储在 Ti中。将这组前缀
称为 Di。




                                                         21
图 8将图 11.7中的固定步长 Trie转换为可变步长 Trie，通过将每个 Trie节点的步长与指向该节点
的指针一起编码。注意，最左边的叶节点现在只包含四个位置（而不是八个），从而将位置数量
从 24减少到 20。
Figure 8Transforming the fixed-stride trie of Fig.11.7 into a variable-stride trie by encoding
the           stride         of           each             trie          node            along
withapointertothenode.Noticethattheleftmostleafnodenowonlycontainsfourlocations(insteado
feight),thusreducing the number of locations from 24 to 20




                       图 9 通过动态规划选择最优的可变步长 Trie。
           Figure9 Pickinganoptimumvariable-stridetrieviadynamicprogramming.

                                                                                           22
假设我们可以递归地找到高度为 h-1和数据库 Di的最优 Ti。在根节点使用一次内存访问后，只剩
下 h-1次内存访问来导航每个子 TrieTi。设 Ci表示最优 Ti所需的存储成本，以数组位置计数。然
后，对于固定的根步长 s，结果最优 Trie的成本 C(s)是 2s（根节点在数组位置中的成本）加上所
有子 TrieTi的成本之和（∑ Ci）。因此，初始步长的最优值是 1⩽ s⩽ 32中使 C(s)最小的值。
递归的简单使用会导致重复的子问题。为避免重复的子问题，算法首先构建一个辅助的单比特
Trie。注意，图 11.9中的任何子 TrieTi必须是单比特 Trie的一个子 TrieN。然后算法使用动态规划
构建从高度 1到 h的每个子 TrieN在原始单比特 Trie中的最优成本和 Trie步长，从最小高度的子
Trie到最大高度的子 Trie自下而上地构建。最终结果是高度为 h的根节点（单比特子 Trie）的最优
步长。详情见 Srinivasan和 Varghese（1999）。
最终，算法的复杂度很容易看出是 O(N*W2*h)，其中 N是原始数据库中前缀的数量，W是目的地
址的宽度，h是期望的最坏情况高度。这是因为在单比特 Trie中有 N*W个子 Trie，每个子 Trie的
高度从 1到 h，每个解决方案需要在最多 W个初始步长 s的可能选择中进行最小化。注意，复杂度
在 N（写作时最大约 150,000）和 h（应较小，最多 8）上是线性的，但在地址宽度（目前 32）
上是二次的。实际上，地址宽度的二次依赖性并不是主要因素。
   例如，Srinivasan和 Varghese（1999）显示，使用高度为 4，优化后的 MAE-East数据库需要
423KB的存储空间，而未优化版本需要2003KB的存储空间。未优化版本使用“自然”步长长度
   8,8,8,8。动态规划在 300MHz的 PentiumPro上运行需要 1.6秒。对于固定步长 Trie，动态程序
更简单，只需 1毫秒即可运行。然而，使用固定步长需要 737KB的存储空间，而不是 423KB。 显
然，1.6 秒对于每次更新运行动态程序来说太长了，无法实现毫秒级更新（Labovitz 等，1997）。
然而，骨干网的不稳定性是由病态情况引起的，在这些情况下，路由器临时被淹没，导致同一组前
缀 S 被重复插入和删除（Labovitz 等，1997）。由于我们无论如何都必须为包括 S 在内的整个集
合分配内存，因此当 S 被删除时，Trie 在内存使用上次优的事实是无关紧要的。另一方面，管理
者添加或删除新前缀的速度似乎更可能是以天为单位。因此，每天运行几秒钟的动态程序是合理
的，并且不会过度影响最坏情况的插入和删除时间，同时仍然允许合理优化的 Trie。



11.5.3 增量更新


                                                            23
图 10 层次压缩（LC）Trie方案将单比特 Trie递归分解为尽可能大的完整子 Trie（左图）。每个
完整子 Trie中的子节点（由虚线框表示）然后被放置在一个 Trie节点中，形成一个特定于所选数
据库的可变步长 Trie。 Figure10 The level-compressed (LC) trie scheme decomposes the 1-
bit trie recursively into full subtries of the largest size
possible(left).Thechildrenineachfullsubtrie(shownbythedottedboxes)arethenplacedinatrienode
toformavariable-stridetrie that is specific to the database chosen.

多比特 Trie有简单的插入和删除算法。考虑添加一个前缀 P。算法首先模拟在新前缀 P的位字符
串上的查找，直到并包括前缀 P中的最后一个完整步长。查找将终止于最后一个（可能不完整的）
步长或到达一个空指针。因此，对于向图 11.7 的数据库添加 P10 = 1100*，查找跟随 110 指针并
在最左边的叶 Trie节点 X处终止。
为了插入和删除的目的，对于多比特 Trie 中的每个节点 X，算法维护一个对应的单比特 Trie，前
缀存储在 X 中。这个辅助结构不需要在高速内存中。此外，对于每个节点数组元素，算法存储其
当前最佳匹配的长度。
在确定必须将 P10添加到节点 X之后，算法将 P10扩展到 X的步长。任何扩展到 P10的数组元素
（当前标记为长度小于 P10的前缀）必须被 P10覆盖。因此，在添加 P10 = 1100*时，算法必须
将 0*的扩展添加到节点 X中。具体来说，节点 X中的 000和 001条目必须更新为 P10。
如果查找在到达前缀中的最后一个步长之前终止，算法会创建新的 Trie 节点。例如，如果添加前
缀 P11 = 1100111*，当在 011条目处找到一个空指针时，查找在节点 X处失败。然后，算法在此位
置创建一个新指针，该指针指向包含 P11的新 Trie节点。然后在这个新节点中扩展 P11。
删除类似于插入。插入和删除的复杂度是执行查找的时间（O(W)）加上完全重建 Trie节点的时间
（O(S)，其中 S是 Trie节点的最大大小）。例如，使用 8位 Trie节点，后者的成本将需要扫描大

                                                                                       24
约 28= 256个 Trie节点条目。因此，为了允许快速更新，关键是还要限制前面描述的动态程序中
任何 Trie 节点的大小。



11.5 层次压缩（LC）Trie




                        图 11LCTrie的数组表示。
                 Figure11 ArrayrepresentationofLCtries.



LCTrie（Nilsson和 Karlsson，1998）是一种可变步长 Trie，其中每个 Trie节点都不包含空条
目。构建 LCTrie的方法是首先找到允许无空条目的最大根步长，然后在子 Trie上递归重复此过
程。图 11.10展示了这个过程，从左侧的单比特 Trie开始，最终形成右侧的 LCTrie。注意，P4和
P5 构成了最大的可能完整根子 Trie——如果根步长为 2，那么前两个数组条目将为空。这样做的
目的是避免空数组元素，以最小化存储。然而，一般的可变步长 Trie 更具可调性，允许在内存和
速度之间进行权衡。
例如，使用 1997年 MAE-East的快照数据，LCTrie的高度为 7，需要 700KB的内存。相比之
下，最优可变步长 Trie（Srinivasan和 Varghese，1999）的高度为 4，使用 400KB的内存。还
记得，最优可变步长计算给定目标高度的最佳 Trie，如果 LC Trie 在其高度上是最优的，它将生
成LC Trie。

                                                          25
在最终形式中，可变步长 LCTrie节点按广度优先顺序布局（首先是根，然后是第二层从左到右的
所有 Trie节点，依此类推），如图 11.11右侧所示。每个指针变成一个数组偏移量。数组布局和完
整子 Trie的要求使得最坏情况下的更新很慢。例如，删除图 11.10中的 P5会导致子 Trie分解的变
化。更糟的是，这几乎会导致图 11.11 中数组表示的每个元素都向上移动。



11.6 Lulea压缩 Trie




           图 12使用 Lulea位图压缩方案压缩图 11.7的根节点（经过叶推之后）。
Figure12
CompressingtherootnodeofFig.11.7(afterleafpushing)usingtheLuleabitmapcompressionschem
e.

虽然 LCTrie和可变步长 Trie尝试通过在每个节点改变步长来压缩多比特 Trie，但这两种方案都有
问题。LCTrie通过使用完整数组避免了由于空数组位置而浪费内存，但这也增加了 Trie的高度，
使其无法调整。另一方面，可变步长 Trie可以通过调整步长来保持较短的高度，但代价是由于
Trie节点中的空数组位置而浪费内存。Lulea方法（Degermark等，1997）是一种多比特 Trie方
案，它使用固定步长的 Trie节点，但通过使用位图压缩显著减少存储。
我们知道，具有重复的字符串（例如，AAAABBAAACCCCC）可以使用表示重复点的位图（即
10001010010000）和压缩序列（即 ABAC）进行压缩。同样，图 11.7 的根节点包含由于扩展而产
生的重复序列（P5, P5, P5, P5）。
Lulea方案（Degermark等，1997）通过使用位图和压缩序列来压缩重复信息，避免了这种明显的
浪费，而不会在查找时间上付出高昂代价。例如，该方案仅使用 160KB内存存储 MAE-East数据
库。这使得整个数据库可以放入昂贵的 SRAM 或片上内存。然而，它在插入时间上付出了高昂的
代价。
                                                                                   26
一些扩展的 Trie 条目（例如，图 11.7 根节点中的 110 条目）有两个值，一个指针和一个前缀。为
了使压缩更容易，算法首先通过将前缀信息下推到 Trie 叶节点，使每个条目只有一个值。由于叶
节点没有指针，我们在叶节点只有下一跳信息，而在非叶节点只有指针。这个过程称为叶推。
例如，为了避免图 11.7 根节点中 110 条目中存储的额外前缀，存储的前缀 P9 被推到最左边的
Trie 节点中的所有条目，除了 010 和 011 条目（它们继续包含 P3）。同样，根节点条目 100 中
的存储前缀 P8 被下推到最右边的 Trie 节点的 100、101、110 和 111 条目中。一旦完成，每个节
点条目只包含一个存储前缀或一个指针，而不是两者兼有。
Lulea方案从一个概念上的叶推扩展Trie开始，用一个值替换连续的相同元素。使用节点位图
（其中 0对应于被移除的位置）来允许对压缩节点进行快速索引。
考虑图 11.7中的根节点。经过叶推后，根节点的序列为 P5,P5,P5,P5,ptr1,P1,ptr2,P2



（ptr1是指向包含 P6和 P7的 Trie节点的指针，ptr2是指向包含 P3的节点的指针）。在将连续的
值替换为第一个值后，我们得到 P5, -, -, -, ptr1,P1,ptr2,P2，如图 11.12中间框所示。最右框显示
了最终结果，其中有一个位图指示移除的位置（10001111）和一个压缩列表（P5,ptr1,P1,ptr2,
P2）。
如果在原始（未扩展的）Trie 节点中有 N 个原始前缀和指针，则压缩节点中的条目数不会超过
2N + 1。直观地说，这是因为 N个前缀将地址空间划分为最多 2N + 1个不相交的子范围，每个子
范围最多需要一个压缩节点条目。
搜索使用步长指定的位数索引到当前 Trie节点，从根节点开始，直到遇到空指针为止。然而，在
跟随指针时，未压缩索引必须映射到压缩节点中的索引。这个映射通过计算节点位图中的位数来完
成。
      考虑图 11.12右侧的数据结构，并搜索一个以 100111开头的地址。如果我们只处理图 11.12左
    侧的未压缩节点，可以使用 100索引到第五个数组元素以获得 ptr1。然而，现在我们必须从图
                                                             11.12
右侧的压缩节点表示中获得相同的信息。
相反，我们使用前三位（100）索引到根节点位图。由于这是设置的第二个位（算法需要计算给定
位之前设置的位数），算法索引到压缩节点的第二个元素。这产生一个指向最右边 Trie 节点的指
针 ptr1。接下来，假设图 11.7最右边的叶节点（经过叶推后）也以相同方式压缩。节点包含序列
P7,
                                                                27
P6,P6,P6,P8,P8,P8,P8。因此，相应的位图是 11001000，压缩序列是 P7,P6,P8。
因此，在最右边的叶节点中，算法使用目标地址的下三位（111）索引到位 8。由于这一位是 0，
搜索终止：在等效的未压缩节点中没有要跟随的指针。然而，为了检索该节点处的最佳匹配前缀
（如果有的话），算法必须找到该条目之前存储的任何前缀。
如果进行扩展，这将是微不足道的，因为值 P8会扩展到 111条目，但由于在压缩版本中 P8值的扩
展序列已被单个 P8值替换，算法需要更多工作。因此，Lulea算法计算位置 8之前设置的位数
（恰好是 3），然后索引到压缩序列的第三个元素。这给出了正确的结果 P8。
Lulea论文（Degermark等，1997）描述了一种使用固定步长 16、8和 8的 Trie。但算法如何高效
地计数大型位图中设置的位数，例如 16 位步长需要使用的 64K 位大小的位图？在继续阅读之
前，尝试使用 P12（为速度添加状态）和 P2a（预计算）原则回答这个问题。
为了加快计数设置的位数，算法为每个位图附加一个摘要数组，该数组包含与位图固定大小块相关
的已设置位的累积计数（预先计算）。使用 64位块，摘要数组占用的存储空间可以忽略不计。计
算位置 i之前设置的位数现在需要两个步骤。首先，访问位置 j的摘要数组，其中 j是包含位 i的
块。然后访问块 j并计算块 j中到位置 i的位数。这两个值的总和就是所需的计数。
虽然 Lulea论文使用 64位块，但图 11.13中的示例使用 8位块。大位图从左到右显示，起始值为
10001001，是从上数第二个数组。每个 8 位块有一个摘要计数，显示为位图上方的一个数组。块
i 的摘要计数计算位图前面块的累积位数（不包括块 i）。
因此，第一个块的计数为 0，第二个块的计数为 3（因为 10001001有三个设置的位），第三个块
的计数为 5（因为 10000001 有两个设置的位，加上前一个块的值 3，得到累积计数 5）。
考虑搜索图 11.13中位置 X之前设置的位数，X可以写作 J011。显然，X属于块 J。算法首先查找摘
要计数数组以检索 numSet[J]。这得出了不包括块 J的累积设置位数。然后算法检索块 J本身
（10011000），并计算块 J前三位的设置位数。由于块 J的前三位是 100，这得出值 1。最后，所



需的总体位计数是numSet[J]+1。




                                                       28
图 13为了在大型位图（例如 64Kbits）中快速计数已设置的位数，位图被划分为多个块，并预先
计算每个块之前已设置的位数摘要计数。
Figure13Toallowfastcountingofthebitsseteveninlargebitmaps(e.g.,64Kbits),thebitmapisdivide
dintochunksanda summary count of the bits set before each chunk precomputed.



块大小的选择是内存大小和速度之间的权衡。将块设为与位图大小相等会使计数非常慢，而将块设
为一位则需要比原始 Trie 节点更多的存储空间。选择 64 位块大小使得摘要数组的大小仅为原始
节点大小的 1/64，但这需要计算 64 位块内设置的位数。计数可以通过软件中的特殊指令和硬件
中的组合逻辑轻松完成。
因此，节点的搜索首先需要索引到摘要表，然后索引到相应的位图块，以计算压缩节点中的偏移
量，最后从压缩节点中检索元素。这每个节点可能需要三次内存引用，这可能会很慢。
最终的 Lulea方案还基于下一跳值压缩条目（具有相同下一跳值的条目即使匹配不同的前缀也可以
被视为相同）。总体而言，Lulea方案具有非常紧凑的存储空间。使用 1997年 MAE-East数据库
的约 40,000条目快照，Lulea论文（Degermark等，1997）报告将整个数据库压缩到约 160KB，
大约每个前缀 32位。
这是一个非常小的数字，因为在数据库中预计每个前缀至少使用一个 20 位指针。紧凑的存储空间
是一个巨大优势，因为它允许前缀数据库潜在地适应有限的片上 SRAM，这是允许前缀查找到
OC-768 速度的关键因素。
尽管存储空间紧凑，Lulea方案有两个缺点。首先，计数位数每个节点至少需要一次额外的内存引
用。其次，叶推使得最坏情况下的插入时间很长。添加到根节点的前缀可能导致信息被推送到成千
上万个叶节点。我们接下来研究的完整树位图方案，通过放弃叶推并使用每个节点的两个位图来克
服这些问题。



11.7 树位图


                                                                                       29
图 14树位图方案通过在每个节点使用两个位图来实现 Lulea的压缩，同时不牺牲快速插入的性
能。第一个位图描述有效指针与空指针，第二个位图描述内部存储的前缀。
Figure14The tree bitmap scheme allows the compression of Lulea without sacrificing fast
insertions by using two bitmaps per node.The first bitmap describes valid versus null
pointers, and the second describes internally stored prefixes.

树位图（TreeBitmap）方案（Eatherton等，2004）的目标是实现与 Lulea方案相同的存储和速
度，同时增加快速插入的目标。虽然我们认为快速插入不如快速查找重要，但它们显然是受欢迎
的。此外，如果处理插入或删除的唯一方法是重建 Lulea压缩的 Trie，那么路由器必须保留其路由
数据库的两个副本，一个用于构建，另一个用于查找。这可能会将每个前缀的存储成本从 32位增
加到 64位，从而使放置整个数据库在片上 SRAM中的芯片支持的前缀数量减半。
为了实现快速插入，从而避免需要两个数据库副本，必须处理 Lulea中的第一个问题，即叶推。当
插入一个小长度的前缀时，叶推可能导致将前缀推到大量叶节点，导致插入变慢。

11.7.1 树位图想法
树位图方案的第一个主要思想是每个 Trie节点有两个位图，一个用于所有内部存储的前缀，另一
个用于外部指针。图 11.14 展示了图 11.12 中根节点的树位图版本。
回顾图 11.12，原始数据库中的前缀 P8 = 100*和 P9 = 110*在左侧图片中缺失，因为它们已被推
到叶节点以容纳两个指针（ptr1指向包含较长前缀的节点，如 P6 = 1000*，ptr2指向包含较长前
缀的节点，如 P3 = 11001*）。这形成了基本的 LuleaTrie节点，其中每个元素包含指针或前缀，
但不能同时包含两者。这允许使用单个位图来压缩 Lulea 节点，如图 11.12 最右侧所示。
相比之下，图 11.14 中的同一个 Trie 节点被分成两个压缩数组，每个数组都有自己的位图。第一
个数组垂直显示，是一个指针数组，包含一个位图，表示存在非空指针的两个位置，以及一个包含
非空指针 ptr1 和 ptr2 的压缩数组。
第二个数组水平显示，是内部前缀数组，包含所有在前 3位内的前缀的列表。用于这个数组的位图
与Lulea编码非常不同，对于每个可能存储在此节点中的前缀都有一个位设置。可能的前缀按
                                                                                          30
字典顺序列出，从 * 开始，然后是 0* 和 1*，接着是长度为 2 的前缀（00*、01*、10*、11*），
最后是长度为 3 的前缀。当相应的前缀出现在 Trie 节点中时，设置位。
因此，在图 11.14 中，在图 11.12 中被推到叶节点的前缀 P8 和 P9 已被恢复，现在对应于内部前
缀位图中的第 12和第 14位。一般来说，对于一个 r位的 Trie节点，有 2r+1− 1个长度为 r或更短的
可能前缀，需要使用一个 (2r+1− 1) 位图。这个方案得名于内部前缀位图以线性化格式表示 Trie：
Trie的每一行从上到下从左到右捕获。
树位图方案的第二个思想是尽量保持 Trie 节点的小型化，以减少给定步长所需的内存访问量。因
此，Trie 节点是固定大小的，仅包含一个指针位图、一个内部前缀位图和子指针。但存储前缀关
联的下一跳信息怎么办？
技巧是将每个 Trie 节点中存储的内部前缀关联的下一跳信息存储在与该 Trie 节点关联的单独数组
中。将下一跳指针放在单独的结果数组中，可能需要每个 Trie节点进行两次内存访问（一次访问
Trie 节点，另一次访问存储前缀的结果节点）。
然而，一个简单的惰性求值策略（P2b）是不在搜索终止前访问结果节点。终止时，算法进行最后
一次访问以获取正确的结果节点。这是对应于路径中最后遇到的包含有效前缀的 Trie 节点的结果
节点。这在末尾只增加了一次内存引用，此外每个 Trie 节点需要一次内存引用。
第三个思想是每个节点只使用一次内存访问，而不像Lulea 需要至少两次内存访问。Lulea 每个节
点需要两次内存访问，因为它使用了较大的步长 8或 16位。这使得位图的大小增加了很多，以至
于唯一可行的计数方式是使用一个必须单独访问的额外块数组。树位图方案通过简单地使用较小步
长的节点（例如 4位）解决了这个问题。这使得位图足够小，以至于整个节点可以通过一次宽访问
来访问（P4a，利用局部性）。可以使用组合逻辑（第二章）来计数位数。

11.7.2 树位图搜索算法
树位图搜索算法从根节点开始，使用目标地址的前 r位（对应于根节点的步长，我们的例子中为
3）在根节点的指针位图中索引到位置 P。如果这个位置有 1，则表示有一个有效的子指针。算法
计算这个 1左边的 1的数量（包括这个 1），并将这个计数记为 I。由于子指针块的起始位置（例
如， y）和每个 Trie节点的大小（例如，S）已知，可以计算子节点的指针为 y + (I*S)。
在移动到子节点之前，算法还必须检查内部位图，看看是否有一个或多个存储的前缀对应于多比特
节点路径中的位置 P。例如，假设 P是 101，并且在根节点位图中使用了 3位步长，如图 11.14所
示。算法首先检查是否有存储的内部前缀 101*。由于 101 对应于内部前缀位图中的第 13 位，算法
                                                      31
可以检查该位置是否有 1（示例中有 1）。如果该位置没有 1，算法将回退检查是否有对应于 10的
内部前缀。最后，如果没有 10前缀，算法检查前缀 1。
这种搜索算法似乎需要多次迭代，与内部位图长度的对数成正比。然而，对于硬件中高达 512位左
右的位图，这只是简单的组合逻辑问题。直观上，这种逻辑执行所有迭代并行，并使用优先编码器
返回最长匹配的存储前缀。
一旦知道在 Trie节点中有一个匹配的存储前缀，算法不会立即从与 Trie节点关联的结果节点中检
索相应的下一跳信息。相反，算法移动到子节点，同时记住存储前缀的位置和相应的父 Trie节
点。目的是记住搜索路径中最后一个包含存储前缀的 Trie 节点 T 及相应的前缀位置。



当搜索遇到扩展位图中相应位置设置为 0 的 Trie 节点时终止。这时，算法最终访问对应于 T 的
结果数组，读取下一跳信息。减少内存访问宽度的进一步技巧在 Eatherton 的硕士论文
（Eatherton, 1995）中有所描述，其中包括许多其他有用的想法。
直观上，树位图中的插入非常类似于没有叶推的简单多比特 Trie 中的插入。前缀插入可能导致
Trie节点完全更改；创建新节点的副本并以原子方式链接到现有 Trie中。Eatherton等（2004）的
压缩结果表明，树位图在压缩和速度方面具有 Lulea方案的所有特点，同时还具有快速插入的能
力。树位图还可以根据硬件实现进行调整，从使用类似 RAMBUS 的内存到片上 SRAM。

11.7.3 PopTrie：一种替代的位图算法
在最初的 Lulea 和 Tree 位图算法提出多年后，一种新的位图算法称为 Poptrie（Asai 和
Ohara， 2015）被提出。该算法得名于人口计数指令（称为 popcount），可以在现代 CPU上计
数设置的位数
（例如在 64 位机器字中）。该算法针对软件实现，可能用于网络功能虚拟化。动机是由于 Tree
位图的存储前缀不是简单的线性位图，因此在软件中难以计算。Poptrie论文（Asai和 Ohara，
2015）描述了一些实验，表明 Tree 位图在软件中相当慢，并提出了一种替代方案。
Poptrie 算法的替代方案是再次使用激进的叶推，因此插入相比 Tree 位图再次变慢。在 Poptrie
中，每个叶节点存储一个关联前缀（对应于其最具体的祖先前缀）。因此，每个多比特 Trie 元素
要么是一个带有关联叶推前缀的空指针，要么包含一个指向后代节点的指针。每个节点都有一个基
本位图，其中的 1表示有效指向后代的指针，0表示存储的前缀。因此，每个 Trie节点都有基本位

                                                    32
图和两个基指针：第一个基指针（base0）指向存储前缀数组的开始，第二个基指针（base1）指
向后代指针数组的开始。
Poptrie 将 IP地址分成 6 位块（一个统一的 6 位步长），所以每个基本位图是 64 位。在搜索期
间，当使用地址的当前块索引位图时，如果该位是 1，则表示有一个后代节点。在这种情况下，搜
索继续通过计数 1的数量并使用计数索引到压缩的指针数组，其起始位置是 base1。如果索引位
是 0，搜索以存储前缀终止。搜索现在通过计数 0 的数量并索引到 base0 数组来检索与最长匹配
关联的下一跳。由于Poptrie 使用 64位位图，计数可以通过popcount 指令在 1个周期内高效地在
软件中完成。
如描述的基本方案由于非常激进的叶推需要大量存储。但许多连续条目具有相同的存储前缀。与我
们对 Lulea的描述类似，可以使用第二个位图（称为 LeafVector）进行一种形式的游程编码。因
此，如果在 base0数组中有一串连续存储的前缀，例如 P4,P4,P4,P5...，则该数组中的 P4副本被
一个 LeafVector位图 1001...替换。
总之，所有三种方案都使用不同但密切相关的位图语义。Lulea 使用单个位图和大步长，但使用摘
要位图来恢复速度。Tree位图不进行叶推，但使用两个位图，一个用于压缩指针，另一个 Tree位
图表示存储的前缀。Poptrie进行激进的叶推，并使用两个位图：一个区分指针和前缀，一个压缩
连续存储的前缀。
虽然 Poptrie论文显示其在软件中的性能优于 Tree位图，但存在两个问题。首先，不清楚 Lulea是
否不能通过较小（例如 6位）步长模拟 Poptrie的转发性能，并使用 popcount指令计数位。其
次，正如我们下面将看到的，当前软件中最快的方法使用不同的策略，基于前缀范围的二分查找。



11.8 基于范围的二分查找




图 15 在一个样本数据库的小子集上进行二分查找，该子集仅包含前缀为 P4 = 1*和 P1 = 101*的
数据。
                                                     33
Figure15
Binarysearchonvaluesofatinysubsetofthesampledatabase,consistingofonlyprefixesP4=1*andP
1=101*.

到目前为止，我们所有的方案（单比特 Trie、扩展 Trie、LC Trie、Lulea Trie、树位图）都是
Trie 变体。是否还有其他算法范式（P15）来解决最长匹配前缀问题？确切匹配是前缀匹配的一种
特殊情况。二分查找和哈希（Cormen 等，1990）是确切匹配的常用技术。因此，我们应该考虑
将这些标准的确切匹配技术推广到前缀匹配。本文部分讨论了二分查找的改编，下一部分将讨论哈
希的改编。
在范围上的二分查找（Lampson等，1998）中，每个前缀表示为一个范围，使用范围的起点和终
点。因此，N个前缀的范围端点将地址空间划分为 2N+1个不相交的区间。该算法（Lampson等，
1998）使用二分查找来找到目标地址所在的区间。由于每个区间对应一个唯一的前缀匹配，算法
预先计算这个映射并将其存储在范围端点中。因此，前缀匹配需要 log2(2N)次内存访问。
考虑一个只有两个前缀的小路由表，P4 = 1*和 P1 = 101*。这是图 11.6中使用的数据库的一个小子
集。图 11.15 显示了如何将二分查找数据结构构建为一个表（左）和一个二叉树（右）。
该方案的起点是将前缀视为地址范围。为简单起见，假设地址是 4位而不是 32位。因此，P4
= 1*范围是 1000到 1111，P1 = 101*范围是 1010到 1011。接下来，在添加整个地址空间的范围
（0000到 1111）之后，所有范围的端点被排序到一个二分查找表中，如图 11.15左侧所示。
在图 11.15 中，范围端点垂直绘制在左侧。图中还显示了每个前缀覆盖的范围。接下来，每个端点
关联两个下一跳条目。最左边的条目称为 > 条目，对应于大于端点但小于下一个范围端点的地
址。最右边的条目称为 = 条目，对应于正好等于端点的地址。
例如，显然任何大于或等于 0000 但小于 1000 的地址都不匹配任何前缀。因此，0000 对应的条
目是-，表示没有下一跳。同样，任何大于或等于 1000但小于 1010的地址必须匹配前缀 P4 =
1*。
唯一微妙的情况是 1011的条目。如果一个地址严格大于 1011但小于下一个条目 1111，则最佳匹配
是 P4。因此，> 指针是 P4。另一方面，如果地址正好等于 1011，则最佳匹配是 P1。因此，=指
针是 P1。
整个数据结构可以构建为一个二分查找表，每个表条目有三个项目，包括一个端点、一个 > 下一
跳指针和一个= 下一跳指针。该表最多有2N个条目，因为每个N个前缀可以插入两个端点。因


                                                                                    34
此，在预先计算下一跳值后，可以使用端点值上的二分查找在 log2 2N 时间内搜索表。或者，表
可以绘制为二叉树，如图 11.15 右侧所示。每个树节点包含端点值和相同的两个下一跳条目。
到目前为止的描述表明，基于值的二分查找可以在 log2 2N 时间内找到最长前缀匹配。然而，可
以使用更高基数的二叉树（如 B树）来减少时间。虽然这样的树需要更宽的内存访问，但对于允
许快速访问连续内存位置的 DRAM来说，这是一个有吸引力的权衡（P4a）。
计算几何（Preparata和 Shamos，1985）提供了一种称为范围树的数据结构，用于找到最窄的范
围。范围树提供快速插入时间和快速 O(log2N)搜索时间。然而，似乎没有简单的方法增加范围树
的基数以获得 O(logMN)搜索时间，其中 M>2。
如前所述，可以使用栈和附加的 Trie轻松在线性时间内构建此数据结构。不难看出，即使使用平
衡二叉树（见练习），添加短前缀也会改变表中大量前缀的 > 指针。Warkhede等（2001）描述
了一种允许快速插入和删除的对数时间技巧。
基于前缀值的二分查找在与多比特 Trie 相比时显得有些慢。然而，与其他受专利限制的 Trie 方案
不同，二分查找没有这种限制。因此，至少有一些厂商将该方案实现到硬件中。在硬件中，使用宽
内存访问（以减少对数的基数）和流水线（允许每次内存访问进行一次查找）可以使该方案足够
快。
1998 年撰写的原始论文（Lampson 等，1998）提出了两个优化以提高软件中的速度。首先，它
建议通过预先计算的最佳匹配前缀表来改进基本二分查找方案的最坏情况内存访问次数。论文显
示，使用数组作为前端的简单技巧将每个分区表中的最大前缀数从 38000 多个减少到 336 个，
从而使二分查找更快（10 次内存访问，相对于原始表大小为 N时的 log2 N + 1 次）。其次，它建
议使用更大的基数而不是二叉树，并利用 Pentium 处理器的缓存行大小使这种 k 路查找在软件中
高效。
最重要的是，Lampson等（1998）提出的基于前缀范围二分查找的高度优化形式，称为DXR
（Zec等，2012），在 2021年首次提出，目前在 2022年使用 8核商品 CPU（AMDR7-1700）上
实现了每秒超过 25 亿次 IP 查找。这使其成为网络功能虚拟化设备的良好构建模块。我们现在描
述 DXR中的新优化。



11.9 带初始查找表的范围二分查找
DXR（Zec 等人，2012 年）中的核心思想仍然是在具有初始表的范围上进行二分查找。因此，如
果根节点有 N 个指针，则可以有指向 N 个不同二分查找表的指针。再次强调，每个前缀都被转换
                                                       35
为一个范围（起始地址和结束地址）。
然而，DXR超越了最初的二分查找范围的想法，使用了多线程和精心的压缩，以巧妙的方式适应
现代 CPU 的高速缓存层次结构。首先，他们使用多线程实现，以便通过增加硬件线程数来隐藏二
分查找的显著成本。其次，他们使用了几种压缩方法来减少内存使用，使得数据库更有可能适合
于 L2/L3缓存。将解析为相同下一跳的相邻地址范围合并。接下来，结束地址可以从下一个区间
的开始派生。因此，每个条目只需要起始地址（2个字节，因为前 16位或更多位已经通过直接查找
解析），以及下一个跳（1或 2个字节用于索引下一个跳表）。对于仅引用 8位下一跳索引并且对
应于
最多 24位前缀长度的块，可以实现进一步的双倍压缩。对于这些块，范围起始地址可以压缩为 8
位
或更少，这是一个很大的节省。最后，他们尝试使用超过 216的初始表大小，并通过使用 16到 21



位之间的位数找到更好的数字。




例如，使用 16位初始表是紧凑的，并且在合成测试中每个前缀的开销小于 2个字节，并且超过
100百万次每秒的查找（Mlps）在单个商品 CPU核心上（DXR，2022）。使用 21位允许在单个
CPU核心上每秒 200Mlps，以增加内存占用的代价，并提供超过 20亿次每秒的总体吞吐量
（DXR，2022）使用 8个核心。FreeBSD中合并的 DXR版本使用了一个两阶段的 trie，在进行二
分查找之前进行了 16-4的分割。



需要注意的是，分区（分成多个前缀搜索表）还降低了更新的成本，因为前缀添加或删除前缀 P
不会影响所有分区表，而只会影响与 P匹配的根条目指向的表。这意味着一个全面构建的成本要比
增量更新要高得多：据报道，对于 16位初始表，全面构建的成本为 70毫秒，对于 20位初始步
                                                     36
长，成本上升到 300毫秒（Zec和 Mikuc，2017年），而增量更新的平均成本据报道为随机测试
（Zec等人，2012 年）平均仅为 10 微秒。



DXR 表现出色的两个根本原因似乎有两个。首先，更大的原因是仔细的压缩使其适应了高速缓存
层次结构；随着表大小的增长，现代处理器的缓存大小也应该增长。其次，二分查找所花费的看似
长时间似乎被线程并行性所隐藏。最近更新的 Zec和 Mikuc（2017）建议，随着 CPU扩展到更多
的硬件线程（例如 36 个线程），这可能会轻松将 DXR 的吞吐量加倍。



尽管DXR在IPv4 方面取得了成功，但不太可能直接适用于IPv6，因为IPv6 具有 128位键和更稀疏
的地址空间。然而，值得指出的是，原始论文 Lampson 等人（1998 年）建议对长标识符进行多
列二分查找（请参阅本章有关多列搜索的练习）。多列搜索与多线程（以及 DXR中使用的压缩思
想）可能是快速软件 IPv6 实现的基础，但需要进行更多的工作。



DXR的代码可在网上找到，并已集成到 FreeBSD中，这使得它可以作为基于 FreeBSD的网络功能
虚拟化设备的一部分使用。



11.10 对前缀长度的二分查找



                                                   37
               图 16 从简单的线性搜索可能的前缀长度到二分查找。
     Figure16 Fromnaivelinearsearchonthepossibleprefixlengthstobinarysearch.

在这一节中，我们将另一种经典的精确匹配方案——哈希，改编成最长前缀匹配。对前缀长度进行
二分查找使用 log2W个哈希，其中 W是最大前缀长度。这可以为 128位 IPv6地址提供非常可扩展
的解决方案。对于 128位前缀，该算法仅需要七次内存访问，而使用具有 8位步长的多位
字典树则需要 16次内存访问。为此，该算法首先根据长度将前缀分开到不同的哈希表中。更具体
地说，它使用一个数组 L 来存储哈希表，使得 L[i] 是指向长度为 i 的所有前缀的哈希表的指针。
假设有相同的小型路由表，仅包含两个前缀 P4 = 1* 和 P1 = 101*，分别为长度为 1 和 3，这是图
11.15中使用的一个小的子集，回顾一下这是图 11.6的一个小的子集。哈希表的数组在图 11.16
的顶部框架（A）水平显示。长度为 1的哈希表存储 P4，垂直显示在左侧，并且在数组中的位置 1



指向它；长度为 3的哈希表存储 P1，显示在右侧，并且在数组中的位置 3指向它；长度为 2的哈
希表为空，因为没有长度为 2的前缀。
                                                                               38
在简单搜索算法中，对于地址 D 的搜索将从最长长度的哈希表 l（即 3）开始，将 D 的前 l 位提
取到 Dl中，然后在长度为 l的哈希表中搜索 Dl。如果搜索成功，则找到了最佳匹配；否则，算法
将考虑下一个较小的长度（即 2）。算法在可能的前缀长度集合中以递减的顺序移动，直到找到匹
配或耗尽长度。
简单的方案有效地在不同的前缀长度中进行线性搜索。类比表明一个更好的算法：二分查找
（P15）。但是，与对前缀范围的二分查找不同，这是对前缀长度的二分查找。差异很大。对于32
个长度，在最坏的情况下，对长度的二分查找需要五个哈希；而对于 32,000个前缀，在前缀范围
上
进行二分查找需要 16次访问。二分查找必须从中位数前缀长度开始，每个哈希必须将可能的前缀
长度分为两半。哈希搜索只给出两个值：找到和未找到。如果在长度为 m的位置找到匹配项，则
必须搜索长度严格大于 m 的匹配项。相应地，如果未找到匹配项，则搜索必须继续在严格小于
m 的长度的前缀中进行。
例如，在图 11.16的部分（A）中，假设搜索从中位数长度为 2的哈希表开始，对于以 101开头
的地址。显然，哈希搜索找不到匹配项。但是在长度为 3的表中有更长的匹配项。由于只有匹配才
会导致搜索移到右半部分，因此必须引入“人为匹配”或标记，以在可能有更长匹配项时强制搜索
移到右半部分。因此，部分（B）在长度为 2的表中引入了一个粗体标记条目 10，对应于前缀
P1=101的前两位。实质上，为了提高速度，添加了状态。这些标记允许在中位数处的探测失败时
排除所有长度大于中位数的长度。对以 101开头的地址 D进行搜索是正确的。在图 11.16的部分
（B）中，假设搜索从中位数长度为 2的哈希表开始，对于以 101开头的地址。在长度为 2的表中
搜索 10（在图 11.16的部分（B），结果是匹配；搜索继续到长度为 3的表，在 P1找到匹配项，并
终止。通常，对于要搜索 P的搜索二分查找将访问的所有长度必须在 P搜索中放置标记。这只增加
了对数数量的标记。对于长度为 32的前缀，只需要在长度为 16、24、28和 30的长度上放置标
记。
不幸的是，该算法仍然不正确。虽然标记可以导致潜在的更长前缀，但它们也可能导致搜索跟随错
误的线索。考虑对在图 11.16 的部分（B）中前三位为 100 的地址 D进行搜索。由于中位数表包
含 10，搜索中间哈希表会导致匹配。这迫使算法在第三个哈希表中搜索 100并失败。但是正确的
最佳匹配前缀在第一个哈希表中，即 P4=1*。标记可能导致搜索陷入死胡同！另一方面，在左半
部分进行回溯搜索将导致线性时间。
为了确保对数时间，每个标记节点 M包含一个变量 M.bmp，其中 M.bmp是与字符串 M匹配的最
长前缀。当 M 插入其哈希表时，这是预先计算的。当算法跟随标记 M 并搜索长度大于 M 的前缀
时，如果算法找不到这样的更长前缀，则答案是 M.bmp。实质上，每个标记


                                                 39
的最佳匹配前缀都是预先计算的。当使用标记获得匹配项时，避免在获得匹配项时搜索中位数以下
的所有长度。包含前缀 P4 和 P1 的数据库的最终版本显示在图 11.16 的部分（C）中。为长度为
2的哈希表添加了指向字符串 10的最佳匹配前缀（即 P4 = 1*）的 bmp字段。因此，当算法搜索
100并在中位数长度为 2的表中找到匹配项时，它会记住搜索长度为 3的表之前遇到的最后一个标
记的 bmp条目 P4。当搜索失败时（在长度为 3的表中），算法返回上一个遇到的标记的 bmp字
段（即 P4）。
从头开始构建简单二分查找数据结构的平凡算法如下。首先确定不同的前缀长度；这确定了要



搜索的长度序列。然后依次将每个前缀 P添加到与 length(P)对应的哈希表中。对于每个前缀，还
要向所有二分搜索将要访问的哈希表（如果尚不存在）中添加一个标记。对于每个这样的标记
M，使用辅助 1 位字典树确定 M 的最佳匹配前缀。Waldvogel 等人（1997 年）描述了进一步的
改进。
虽然在最坏的情况下，搜索算法对IPv4需要五次哈希表查找，但我们注意到，在预期的情况下，大
多数查找应该需要两次内存访问。这是因为预期情况观察 O1显示，大多数前缀要么是 16位，要么
是 24 位（至少在今天）。因此，对 16 位进行二分搜索，然后对 24 位进行二分搜索就足够了。
哈希的使用使得对前缀长度进行二分查找在硬件中有些难以实现。然而，它在大前缀长度（例如
IPv6 地址）方面的可扩展性是值得注意的。



11.11 带硬件辅助的前缀长度线性搜索
上一节表明，我们可以通过添加标记和其他机制使前缀长度的二分搜索起作用。不幸的是，这增加
了复杂性并使插入变慢。本节的想法是重新审视简单的前缀长度线性搜索（我们在上一节中放弃
的），但添加一些硬件辅助使其变得实用。
假设设置是具有芯片内存的硬件（例如 1纳秒访问时间）和大量较慢的 DRAM（例如 50纳秒访问
时间）。我们如何能够通过少量（例如 1或 2个）顺序 DRAM访问进行 IP查找？我们甚至可以承
受高达 32次芯片内存访问，因为芯片访问远远小于单个 DRAM访问。我们将描述两种这样的方
案，都基于使用位图来表示每个长度的存储前缀。由于这些位图是紧凑的，它们可以存储在芯片内
存中；对于每个长度，使用诸如 d-left（Broder和 Mitzenmacher，2001年）这样的方案，在芯
片外的 DRAM中存储相应的哈希表（在第 10.3.3节中描述）。例如，如果只有两个前缀 P4=1∗和
                                                       40
P1=100∗。长度为 1的位图将是 10（对应于 P4的第一个位设置），长度为 2的位图将是 0000，
长度为 3的位图将是 00010000（对应于 P1的第四位设置）。
因此，骨架想法是在芯片上进行顺序或并行搜索以找到最长匹配长度，假设为 L。然后，对芯片外
的 DRAM哈希表进行访问，该哈希表存储所有 L位前缀，使用地址的前 L位作为检索下一跳的关
键。显然，索引位图的朴素使用效果不佳，因为在前缀长度为 32时，相应的位图大小为 232，太
大而无法放入芯片内存中。我们将描述两种方案，展示如何处理由更大长度位图引起的这种内存爆
炸。

11.11.1 使用 Bloom过滤器来压缩前缀位图
观察到，特别是在较大的长度上，位图会非常稀疏。因此，每个位图可以通过所谓的 Bloom 过滤
器（Bloom，1970 年）进行压缩。在 Bloom 过滤器中，任何条目 x 都会被哈希小的恒定次数
（通常为 3次）在独立的哈希函数中（比如 h1(x)=i，h2(x)=j，h3(x)=k），并将位 i，j，k设置在
位图中。因此，如果在长度 L处有 NL个前缀，则理论上显示，NL个前缀需要 cNL位，其中 c是一
个小的常数。
虽然我们将在本书后面研究 Bloom 过滤器，当我们调查测量和安全技术时，现在值得注意的是，
Bloom过滤器没有误负，但可能会有误正。当进行查找不在位的条目 y时，可能发生这种情况，对
应的索引 h1(y)，h2(y)，h3(y)都被设置，导致搜索错误地得出 y存在于该 Bloom过滤器中的结
论。



然而，如果 c足够大，误正的概率就很低。
Dharmapurikar等人（2003年）的一篇论文建议使用 32个 Bloom过滤器在芯片上压缩存储的每
个长度上存储的所有前缀，然后进行芯片外查找。如果 i 是最长长度的 Bloom 过滤器匹配，则进
行芯片外访问到芯片外哈希表 i。然而，Bloom 过滤器很少会出现误正，如果出现这种情况，芯片
将尝试匹配的下一个最小 Bloom过滤器长度，依此类推。尽管在最坏的情况下可能会有几次
DRAM访问（这种病态的 IP地址可以被缓存），但预期的 DRAM访问次数接近 1。然而，这种缺
乏确定性性能导致了下一个被称为 SAIL的提案。

11.11.2 SAIL：在枢轴级别之前使用未压缩的位图
SAIL：在枢轴级别之前使用未压缩的位图
                                                        41
SAIL（Yang 等人，2014 年）的起点在于观察到，直到某个长度（他们称之为枢轴长度），可以
轻松地存储所有位图而无需进一步压缩。特别是，他们建议存储所有位图直到 24位只需要
32Mbits。
因此，他们的实现在长度 24 处进行枢轴操作，但这可以更改。虽然 4 M字节很大，但在芯片上
的 SRAM中是可行的。更好的是，所需的芯片内存量保持不变，即使 IPv4数据库任意增大。他们
还观察到，长度大于 24 位的前缀很少：因此，他们使用前缀扩展来将所有这样的前缀存储在
256 元素的多位字典树节点中，这些节点由芯片外的 DRAM中的它们的 24位前缀索引。请注意，
对于这些罕见的情况，我们需要 256个元素的字典树节点，因为 24位前缀的扩展会产生最多 256
个 32位条目。
由于 SAIL还设计为在软件中运行，所以 SAIL从长度 24开始查找。SAIL首先确定最长匹配是否正
好为 24位，少于 24位，或可能大于 24位。如果匹配为 24位，则 SAIL通过在 DRAM中
查找 24 位下一跳表来终止。如果匹配小于 24 位，则从 23 到 1 依次顺序搜索，尝试每个位图直
到找到最长匹配长度，假设为 W；然后在 DRAM中查找 W长度哈希表。最后，如果匹配大于 24
位
（这应该很少发生），SAIL将使用 IP地址的最后 8位来索引相应的芯片外 DRAM中与地址的前
24
位相对应的多位字典树节点。
然后，关键机制是高效地实现一个谓词，以确定最长匹配是否等于 24、小于 24或可能大于 24。
这比看起来要棘手得多。考虑对应于枢轴长度 24的概念性单比特字典树级别。让正在查找的地址
的前 24 位前缀为 P。
有四种情况需要考虑。如果没有与路径 P对应的字典树节点，则最长匹配必须小于 24。其次，如
果有一个字典树节点，但它是一个纯指针，则搜索必须继续在大于 24 的级别上。但是第 2 种情
况可能会失败，因为不保证在更长的长度上找到匹配（就像在前缀长度二分搜索中一样）。第三种
情况是相应的字典树节点是一个存储的前缀但不是一个指针；但在这种情况下，搜索以 24 位匹配
结束。第四种情况是相应的字典树节点既是存储的前缀又是指针。
例如，考虑我们早期数据库的较小版本，其中 P4 = 1∗，P2 = 101∗，P3 = 11001∗，P7 =
1000000∗和新前缀 P10 = 100∗。如果我们假设枢轴级别为 3，当地址的前 3 位是 000 时，发生
第 1 种情况，而当前 3 位匹配时发生第 2 种情况 110（指向 P3 的纯指针）。此外，当匹配的前
3 位为 101 时（它们匹配 P2而不再是前缀）发生第 3种情况。最后，当匹配的前 3位为 100时
（它们匹配 P10并且可能是更长的匹配 P7）发生第 4 种情况。


                                                       42
  在枢轴级别，单个位图无法区分这四种情况，因为单个索引位只有 2种可能的值。因此，SAIL
 使用了一个称为“枢轴推进”的想法，将四种情况减少到三种情况。例如，第 4 种情况（应该很
                                            少
见）可以通过将存储的前缀扩展到长度 25来消除；因此，前缀可以存储在芯片外的多位字典




11.12 在压缩方案中的内存分配

除了二分搜索和固定步幅的多位字典树之外，本章描述的许多方案需要以不同大小分配内存。因
此，如果一个压缩字典树节点从两个增长到三个内存字，插入算法必须释放大小为 2的旧节点并分
配大小为 3的新节点。操作系统中的内存分配有点启发式，使用诸如最佳适配和最差适配等算法，
这些算法不能保证最坏情况的性能。



事实上，所有标准的内存分配器都可能存在非常糟糕的最坏情况碎片化比。分配和释放可能共同导
致将内存分割成一块块的空洞和小的已分配块。具体而言，如果 Max是最大的内存分配请求的大
小，则当所有空洞的大小为 Max-1且所有已分配的块的大小为 1时，最坏情况发生。这可以通过
使用大小为 的请求分配所有内存，然后进行适当的释放来实现。最终结果是只有 1 的内存被保
                                                  Max
证被使用，因为所有未来的请求可能都是大小为Max。
内存分配器的内存使用直接转化为芯片可以支持的最大前缀数量。假设在不考虑分配器的情况下，
可以证明 20MB的芯片上内存可以在最坏情况下支持 640,000个前缀。如果考虑分配器，并且
Max = 32，则芯片只能保证支持 20,000 个前缀！由于 CAM 供应商在编写时正在宣传最坏情况
                                                   43
下 100,000 个前缀的数字，搜索时间为 10 纳秒，更新时间为微秒。因此，考虑到前缀查找的算法
解决方案通常会压缩数据结构以适应 SRAM中，算法解决方案还必须设计内存分配器，使其快速
且最小化地碎片化内存。
有一个古老的结果（Robson，1974 年）表明，没有对内存进行压缩的分配器的利用率不能比

lo
g
M
a
x
 1 更好。例如，对于 Max = 32，这是 20％。由于这是不可接受的，涉及压缩数据结构的算法
2
解决方案必须使用压缩。压缩意味着移动已分配的块以增加空洞的大小。

压缩在操作系统中几乎不会用到，原因如下。如果您移动内存 M的一部分，则必须纠正所有指向
M 的指针。幸运的是，大多数查找结构都是树形结构，其中任何节点最多只被一个其他节点指
向。通过为每个树节点维护父指针，当 M被重新定位时，指向树节点 M的节点可以得到适当的更
正。幸运的是，只有更新时需要父指针，而不需要搜索。因此，父指针可以存储在用于路由处理器
中更新的数据库的芯片外副本中，而不会消耗宝贵的芯片上 SRAM。
即使解决了这个问题，我们仍需要一个简单的算法来决定何时压缩内存的一部分。有关压缩的现有
文献是在垃圾回收的背景下（例如，Refs.Wilson,1992;LaiandBaker,1996）编写的，并且倾向于
使用全局压缩器，在一个压缩周期中扫描所有内存。为了限制插入时间，需要某种形式的局部压缩
器，它只压缩受更新影响的区域周围的少量内存。



11.12.1 基于帧的压缩



                                                         44
图 17 一个查找芯片的模型，在硬件中使用一个通用的 SRAM进行搜索，这个 SRAM可以是片上
或片外的。在某些情况下，外部存储是廉价的低延迟 DRAM。
Figure17Model of a lookup chip that does a search in hardware using a common SRAM
that could be on or off chip.In some cases, the external memory is cheap low-latency
DRAM

为了展示局部压缩方案的简单性，我们首先描述一个极其简单的方案，它进行最小化的压缩，同时
实现了 50% 的最坏情况内存利用率。然后我们将这个方案扩展，将利用率提高到接近 100
在帧合并中，假设所有M 个内存字被分成大小为Max 的M/Max 个帧。帧合并旨在至少保持内存
利用率为 50%。为了做到这一点，所有非空帧应至少 50% 满。帧合并保持以下简单的不变性：
除了一个未填充的帧外，所有帧都至少 50%满。如果是这样，并且M/Max 远大于 1，那么这将产
生接近 50
分配和释放请求（Sikka和 Varghese，2000年）通过向每个单词添加标签来处理，这些标签有助
于识别空闲内存和已分配的块。唯一的额外限制是所有的空洞必须包含在一个帧内；空洞不能跨越
帧。
如果一个帧非空但利用率低于 50%，则称为有缺陷的帧。为了保持不变性，帧合并有一个额外的
指针来跟踪当前有缺陷的帧，如果有的话。现在，如果一个分配导致之前的空帧成为有缺陷的帧，
如果分配小于Max/2，则可以。同样，如果一个释放导致一个填充了超过50%的帧变成不足50%
满，也是如此。例如，考虑一个包含大小为 1 和大小为 Max-1 的两个已分配块的帧，因此利用率
为 100%。如果大小为Max-1的块被释放，利用率可能降低到1/Max。这可能导致两个帧变为有缺
陷的，这将违反不变性。
维护不变性的一个简单技巧是：假设已经有一个有缺陷的帧 F，并且一个新的有缺陷的帧 F’出现
了。通过将 F和 F’的内容合并到 F中来维持不变性。这显然是可能的，因为 F和 F’都不到一
                                                                                   45
半满。请注意，唯一的压缩是局部的，仅限于两个有缺陷的帧 F和 F’。这种局部压缩导致快速的
更新时间。
通过增加帧大小到 kMax并将有缺陷的帧的定义更改为利用率低于 k/(k+1)，可以改进帧合并的最
坏情况利用率。前面描述的方案是 k=1的特例。增加 k会提高利用率，但会增加压缩的成本。在
Sikka和 Varghese（2000年）中描述了更复杂的分配器，其性能更好。



11.13 固定功能查找芯片模型
在Terabit速度下，查找方案经典上是在芯片上而不是在网络处理器上实现的。我们首先在图
11.17 中描述了一个模型，该模型涉及执行搜索和更新的查找芯片，然后描述了一个针对可编程处
理器的替代模型，例如英特尔（从 Barefoot 收购的）Tofino-3（英特尔公司，2022 年）。查
找芯片
（图 11.17）有一个搜索和一个更新过程，两者都访问一个通用的内存，该内存可以是片上的或片
外的（或两者兼有）。更新过程允许增量更新，并且（可能）对每次更新进行内存分配/释放以及
少量的局部压缩。
实际的更新可以完全在芯片上执行，也可以部分或完全在软件中执行。如果半导体公司希望销售一
个使用复杂更新算法（例如，用于压缩方案）的查找芯片，最好也提供一个硬件更新算法。然而，
如果查找芯片是转发引擎的一部分，完全将更新过程委托给线卡上的单独 CPU 可能更简单。
外部存储器可以是 SRAM，就像图 11.17 中所示，但如今更有可能是更便宜的 DRAM。如果需
要，每次访问外部存储器可以相当宽，甚至高达 1000位。使用宽总线这在今天是完全可行的。搜
索和更新逻辑可以在一个内存周期内轻松并行处理 1000位。请记住，宽字访问可以帮助减少搜索
时间，例如在树位图和基于值的二进制搜索方案中。搜索和更新使用时间多路复用来共享存储查找
数据库的公共内存。因此，搜索过程被允许连续进行 S次内存访问，然后更新过程被允许进行 K次
内存访问。如果 S 为 20 且 K 为 2，则这允许更新从搜索中窃取一些周期，同时仅稍微降低搜索
吞吐量。请注意，这会增加搜索的延迟，最坏情况下增加了 K次内存访问。然而，由于搜索过程
很可能是流水线的，因此这可以被视为额外的流水线延迟很小。


                                                   46
芯片有引脚用于接收搜索（例如键）和更新（例如，更新类型、键、结果）的输入，并且可以返回
搜索输出（例如，结果）。该模型可以被实例化用于各种类型的查找，包括 IP 查找（例如，32
位 IP地址作为键和下一跳作为结果）、桥接查找（48位 MAC地址作为键和输出端口作为结果）
和分类（例如，数据包头作为键和匹配规则作为结果）。
每次添加或删除键都可能导致调用释放一个块并分配一个不同大小的块。每个分配请求可以在1
到 Max个内存字的任何范围内。可以分配的内存总量为 M个字。实际的内存可以是片外的、片上
的或两者兼有。显然，即使是片外的解决方案也将在芯片上缓存任何查找树的第一级。片上内存具
有速度快和成本低的优点。不幸的是，由于当前工艺的限制，片上内存在不使用外部内存（无论
是 DRAM 还是 SRAM）的情况下很难支持 100 万个前缀的数据库。
在内部，芯片通常被大量分段。查找数据结构被分成几个部分，每个部分由不同阶段的逻辑并发处
理。因此，内部 SRAM 可能会被分成几个较小的 SRAM，每个流水线级别可以独立访问。
静态地将SRAM在流水线阶段之间分区存在一个问题（Sikka和Varghese，2000年），因为随



着前缀的插入和删除，每个阶段的内存需求可能会变化。一个可能的解决方案是将单个 SRAM分
成相当多的较小 SRAM，可以通过部分交叉开关动态地分配给各个阶段。然而，设计这样一个在
非常高的速度下使用的交叉开关并不容易，尽管这种方法曾在 Procket 路由器中使用过（Chung
等，2004年）。
本章中描述的所有方案都可以进行流水线处理，因为它们都基于树。树中相同深度的所有节点可以
分配到同一流水线阶段。



11.14 可编程查找芯片和 P4
许多路由器使用类似于上一节中的固定功能查找芯片模型，其中固件中可以进行有限的可编程性。
例如，对于广域网 Arista7500R3/7800R3和 7280R3路由器，其文档
（AristaCorporation,2010）指出，“内部的FlexRoute使用算法方法执行查找”。
相比之下，近年来出现了一些芯片，例如英特尔的 Tofino-3（IntelCorporation,2022），其具有
相当大量的 TCAM，并且可以使用高级语言（如 P4）进行编程。在本节中，我们将简要探讨 P4
和 CAM 的 使 用 。 例 如 ， 基 于 英 特 尔 Tofino3.2T 交 换 硅 的 英 特 尔 Aurora710 声 称

                                                                  47
（IntelCorporation, 2022）可以让数据中心将 IP路由表大小增加到 120万条。此外，TCAM分布
在一组可以使用 P4编程的物理阶段之间。
Tofino-3是所谓的可重构匹配表（RMT）（Bosshart等人，2013）方法的一个示例，用于可编程
网络处理器。简而言之，RMT 方法是上一节描述的固定功能芯片的泛化，它在内部是流水线化
的，每个阶段都可以访问片上SRAM，并且专用于单个功能。相比之下，在RMT 架构中，芯片内
部有大量（例如 32 个）匿名的阶段（不专用于任何功能），并且可以编程（可以编程执行数据包
头部的基本功能）。此外，每个阶段都有RAM 和CAM。事实上，例如Tofino-3中的CAM 是如此
丰富，以至于可以支持大量的路由，而无需任何进一步的算法方法。
 当数据包流经 RMT芯片时，每个数据包头部都通过各个阶段进行流式处理，随后的数据包头部也
以同步的方式进行流式处理，以保持流水线的充满。虽然 SRAM和 CAM页面被分配给物理阶段，
但单个逻辑处理阶段（例如树的处理级别）可以通过被分配更多的物理阶段获得更多的内存。每个
        物理阶段不仅可以由内部专家在固件中进行编程，而且还可以由网络运营商使用一种称为
P4 的更高级别语言进行编程（尽管我们的经验是P4 编程也有些神秘）。这种现场可编程性类似
于由可编程门阵列（FPGA）提供的可编程性。但与 FPGA 不同，像 Tofino-3 这样的芯片速度更
快。然而，它们通过提供有限的可编程性，并使用我们现在探讨的 P4 语言来获得速度。

11.14.1 P4开源编程语言
P4（P4开源编程语言）是一种用于“编程与协议无关的数据包处理器”的高级语言。P4的灵活性和
可扩展性改变了以高速度实现网络功能的方式。P4提供了以下功能：
� 可重构性即使 P4程序被部署到可重构的硬件路由器上，操作员仍可以根据需要在现场修改程
序，这与传统的固定功能 ASIC 不同。



� 协议独立性 P4 语言可以指定处理数据包头的方式，而不受所使用的协议的限制。特别是，虽
然 P4可以支持 IP头，但并不限于 IP。因此，它通过允许自定义报头而不仅仅是可编程路由来概括
了 SDN 方法。
� 目标独立性同一个 P4 程序可以部署到各种硬件上，如 CPU、FPGA 和可编程处理器，例
如
Tofino-3。

                                                         48
首先，P4允许使用可编程解析器定义新的报头。例如，如果一个组织想要在以太网和 IPv4报头之
间添加一个新的安全标签（例如 STAG），P4允许操作员定义 STAG的长度和字段，并将其放置
在以太网和 IPv4 报头之间。
 其次，P4定义了一种处理新报头的方式，其中每个报头由一个表进行处理，这些表可以通过精确
匹配或通配位进行索引（从而在精确匹配、最长匹配前缀和 ACL 查找之间进行抽象）。在一个运
行的示例中，STAG字段可以通过一个精确匹配表进行处理，该表会丢弃在表中的某些 STAG值。
  最后，P4 需要指定一个控制流图，将表的处理链接起来，描述基于其报头字段的数据包的整体
处理。因此，在上面的示例中，控制流图将指定在以太网之后由 STAG表处理报头，之后将报头发
送到 IPv4 查找表。
总之，一个 P4程序有三个主要部分：用于解析的报头规范、用于处理报头的表的规范，以及描
述“主”程序的控制流图，该程序指定了控制在表之间的流动方式。
使 P4具有吸引力的是我们之前提到的 RMT（Bosshart等人，2013）架构，该架构提供了一组通
用的匹配-动作阶段，这些阶段排列在线性流水线中。与传统的 FPGA 互连相比，阶段之间的连线
也较短。每个阶段都配备了足够的 CAM 和 RAM 以及多个并行处理器。然后，程序员或编译器可
以将多个物理流水线阶段（甚至是一部分阶段）分配给在控制流图中定义的逻辑表。
需要注意的是，P4仅提供有限的可重新配置性。一些不能通过P4编程的示例包括队列调度
（Tofino-3确实提供了一组队列调度）、有状态处理（例如 NAT）和内容处理（例如入侵检测、
检测内容签名）。最近出现了大量研究，以解决其中的一些问题。例如，数据包事务（Sivaraman
等人， 2016b）尝试允许可编程的有状态处理，而像 PIFO（Sivaraman等人，2016a）这样的方
案则尝试允许可编程的队列调度。我们将在有关数据包调度的章节中考虑这些方案。
尽管存在这些限制，P4似乎很受欢迎。一些重要的用例包括：
首先，简单地重新分配资源可能是有用的。在核心路由器中，需要更多的转发表，而在边缘路由器
中则需要更少的 ACL；而今天的供应商为每个产品制造了不同的硬件，而单个 P4路由器可以简单
地编程为为不同的市场部署提供特定的路由器。
其次，添加新的报头通常非常耗时。例如，VXLAN（Mahalingam 等人，2020）花费了数年时间
才被添加到定制 ASIC 中；但是在 P4 路由器上可以在几小时内完成。
最后，还有一些新的应用程序，例如测量和安全性，可以以企业特定的方式进行编程。我们将在有
关测量和安全性的章节中考虑其中一些问题。

11.14.2 在 P4模型中进行 IP查找

                                                    49
乍一看，像 Tofino-3这样使用 RMT模型的芯片似乎可以轻松实现 IP查找，因为它们拥有大量的
CAM。然而，这些芯片中的 CAM 数量无法扩展到广域数据库。我们是否可以通过将算法方



法与 TCAM 结合来实现更好的效果呢？
例如，MashUp（Rios和 Varghese，2022）是一种算法和硬件技术的“混合”，它使用 TCAM和
SRAM块的树。虽然 CAM树曾被用于降低功耗或更新成本（例如，CoolCAM（Zane等，2003）
和 TreeCAM（Vamanan和 Vijaykumar，2011），但 MashUp的焦点却在于减少 TCAM位数。它
还可以轻松地在现代可重构流水线芯片（如 Tofino-3（Intel公司））中实现。因此，MashUp可
以将像 Tofino-3这样的芯片的覆盖范围扩展到主干 IPv4和 IPv6数据库，或者比今天使用单个逻
辑 TCAM更大的数据中心。Mashup是基于一种称为“平铺树”的技术。换句话说，Mashup考虑了
构建查找树时 TCAM的内部颗粒大小。请注意，TCAM是以单位为单位存在的，具有 W（宽度）
×D
（深度）的颗粒大小，就像内存页面一样。例如，Tofino-3的 W = 44，D = 512。Tofino-3中的
每个单元 TCAM或块可以在单个时钟周期内对高达 44位（其中任何一个可以是通配符 *）的高达
512个前缀进行最长匹配。例如，要适应当前的 IPv6数据库大小（高达 64位的 150,000个前
缀），直接的方法是水平拼接 2个 TCAM块（64<88），然后垂直拼接 300对这些对
（150,00/512约为 300）。
相比之下，MashUp从一个根开始构建 TCAM块的树，该根可以（换句话说是“平铺”）适应TCAM
 块。这将将前缀的逐位分支树（trie）分解为我们可以递归平铺的子树。为了选择在哪个树级别上
定义根 TCAM，观察树中是否有指针，这些指针代表单个逻辑 CAM中不存在的“开销”。第一个想
           法是通过在下游指针数量较少的高度（称为瘦长度）处切割 trie来减少指针开销。例
如，论文（Rios和 Varghese，2022）显示，对于公共广域 IPv6数据库，指针开销仅为总数据库
大小的 0.41%，在高度 20 处上升到 6.84% 在高度 32 处。
第二个想法是通过在单个 TCAM块中打包多达 D个子树来减少瓦片子树的浪费空间。由于剩余的
前缀位可能在不同的子树中重复，因此需要添加一个 log 2 D 位的消歧标记。尽管存在额外的成
本，但打包确保了在最后一个块中的任何浪费空间（可以是一个单元 TCAM）被分摊到 D个子树
上。换句话说，Mashup将子树打包到单个 TCAM块中以减少内部碎片。
第三个想法是进行“货币兑换”，其中“几乎满”子树被 SRAM页面替换，这一过程作者称之
为“RAM混合化”。由于 SRAM无法进行变长匹配，因此必须通过将变长前缀“扩展”到子树中的最

                                                           50
大长度前缀来加以补救。混合化允许进行货币“套利”，因为 SRAM页面比 TCAM块便宜且更加丰
富（在 Tofino-3中是 3比 1）。
Mashup论文展示了使用所有三种技术的 900k前缀的 IPv4数据库大小的结果：瘦长度、标记聚合
和 RAM混合化。一个
具有 16-4-4-8步幅的四阶段树将 TCAM位数减少了 7倍，与使用单个逻辑 CAM的直接解决方案
相比，RAM页面的成本约为 1000个，这远远少于所有 RAM解决方案所需的 RAM。有关更多详细
信息，请参阅 Mashup 论文（Rios 和 Varghese，2022）。



11.15 总结
在介绍了本章中大量孤立的查找变体之后，获得一些透视是很重要的。因此，我们总结一下查找技
术的现状，并调查了设计中使用的常见原则。
查找技术的现状：随着表格大小（高达 100万个前缀）和速度（几 Tbps的聚合吞吐量）的不断



提高，核心路由器中的查找方案面临着严重的压力。MPLS曾经被认为是绕过查找的一种方法，但
现在主要用于避免流量工程的包分类。CAM（内容寻址存储器）甚至正在侵蚀核心路由器领域，
例如 Tofino-3这样的芯片，但仍无法扩展到广域数据库的大小（尽管像 MashupRios和
Varghese（2022）这样的技术可以解决这个问题）。因此，许多核心路由器供应商，如 Arista
（AristaCorporation，2010），仍然使用和设计基于 RAM的算法方案进行查找。
对于核心中速度为 Terabit的没有 CAM的硬件方案，即使经过了这么多年，TreeBitmap仍然似乎
是适当的选择。然而，可能存在限制 Treebitmap使用的专利。这可能会使其他位图方案（例如
Poptrie（Asai和 Ohara，2015）变得吸引人。Treebitmap曾被用于思科的 CRS-1路由器。另一
方面，基于简化的 SAIL（Yang等，2014）的想法可能适用于使用廉价且丰富的片外 DRAM的硬
件设置。
对于数据中心和企业网络来说，在像 Tofino-3这样的芯片中似乎有足够的 CAM，特别是当利用
（例如 Mashup（Rios 和 Varghese，2022）等方案）来降低功耗和 TCAM 位时。
对于软件实现，DXR方案（Zec等，2012）似乎是当今最好的方法，特别是在网络功能虚拟化

                                                          51
（NFV）的背景下。底层的二进制搜索区间方案允许高效的多线程，其吞吐量随着核心数量的增加
而提高，并适合于 L1 缓存。它也没有受到专利的限制。
最后，基于前缀长度的二分查找是吸引人的，因为它具有向大地址长度扩展的属性。不幸的是，它
的哈希使用使得很难保证查找时间。然而，一些厂商在软件实现中使用它。随着 IPv6 的普及，它
可能会成为未来的竞争对手。
最重要的是，如果压缩字典方案每个前缀使用少于 32位，那么与 CAM相比，压缩字典可以使用
更少的晶体管和功耗。这是因为在 CAM中，查找逻辑分布在 N个存储单元中的每个单元中，而在
算法解决方案中，尽管更为复杂，查找逻辑分布在一个小的、恒定数量的阶段中。对这两种方法进
行仔细的 VLSI 缩放分析将非常有用。
基本原则：尽管这是一章关于查找的内容，并且思考查找需要关注当前的市场趋势，但不要忘记这
是一本关于基本原则的书。有可能，未来的路由器可能会使用全光开关和全光处理，甚至用于查
找。在这种情况下，本章描述的具体算法可能会被丢弃；但也许底层的设计原则将会保留。
本章中描述的所有方案都从分而治之的算法原则开始（通过地址、地址范围或前缀长度进行划
分），但通过其他原则获得了效率。首先，大多数方案使用预计算，这种方法以较慢的插入/删除
时间换取较快的搜索时间。这些方案还利用硬件特性，例如宽存储器，利用快速和慢速存储器，交
换内存以节省时间，并优化给定设计的自由度。表 11.1总结了其中一些方案以及其中使用的原则。
其中许多方案还使用了我们可以称之为信息保留转换。例如，用文本字符串替换一系列单向分支，
或将前缀表示为范围的开始和结束。
最后，本章不能对在学术和专利文献中发表的所有有趣的 IP 查找方案进行充分说明。你可以查阅
相关资料。




                                            52
